\documentclass{article}

% Widen the main text block a bit (default article margins are quite large).
\usepackage[letterpaper,margin=1in]{geometry}
\usepackage{parskip}
% if you need to pass options to natbib, use, e.g.:
%     \PassOptionsToPackage{numbers, compress}{natbib}
% before loading neurips_2026

% The authors should use one of these tracks.
% Before accepting by the NeurIPS conference, select one of the options below.
% 0. "default" for submission
% \usepackage{neurips_2026}
% the "default" option is equal to the "main" option, which is used for the Main Track with double-blind reviewing.
% 1. "main" option is used for the Main Track
%  \usepackage[main]{neurips_2026}
% 2. "position" option is used for the Position Paper Track
%  \usepackage[position]{neurips_2026}
% 3. "eandd" option is used for the Evaluations & Datasets Track
 % \usepackage[eandd]{neurips_2026}
 % if you want to opt-in for a de-anomymized submission (not recommended, but OK):
 %\usepackage[eandd, nonanonymous]{neurips_2026}
% 4. "creativeai" option is used for the Creative AI Track
%  \usepackage[creativeai]{neurips_2026}
% 5. "sglblindworkshop" option is used for the Workshop with single-blind reviewing
 % \usepackage[sglblindworkshop]{neurips_2026}
% 6. "dblblindworkshop" option is used for the Workshop with double-blind reviewing
%  \usepackage[dblblindworkshop]{neurips_2026}

% After being accepted, the authors should add "final" behind the track to compile a camera-ready version.
% 1. Main Track
 % \usepackage[main, final]{neurips_2026}
% 2. Position Paper Track
%  \usepackage[position, final]{neurips_2026}
% 3. Evaluations & Datasets Track
 % \usepackage[eandd, final]{neurips_2026}
% 4. Creative AI Track
%  \usepackage[creativeai, final]{neurips_2026}
% 5. Workshop with single-blind reviewing
%  \usepackage[sglblindworkshop, final]{neurips_2026}
% 6. Workshop with double-blind reviewing
%  \usepackage[dblblindworkshop, final]{neurips_2026}
% Note. For the workshop paper template, both \title{} and \workshoptitle{} are required, with the former indicating the paper title shown in the title and the latter indicating the workshop title displayed in the footnote.
% For workshops (5., 6.), the authors should add the name of the workshop, "\workshoptitle" command is used to set the workshop title.
% \workshoptitle{WORKSHOP TITLE}

% "preprint" option is used for arXiv or other preprint submissions
 % \usepackage[preprint]{neurips_2026}

% to avoid loading the natbib package, add option nonatbib:
%    \usepackage[nonatbib]{neurips_2026}

\usepackage[utf8]{inputenc} % allow utf-8 input
\usepackage[T1]{fontenc}    % use 8-bit T1 fonts
\usepackage{hyperref}       % hyperlinks
\usepackage{url}            % simple URL typesetting
\usepackage{booktabs}       % professional-quality tables
\usepackage{amsmath}        % math environments and operators (e.g., \arg\min)
\usepackage{amssymb}        % additional math symbols (e.g., \gtrsim)
\usepackage{amsfonts}       % blackboard math symbols
\usepackage{nicefrac}       % compact symbols for 1/2, etc.
\usepackage{microtype}      % microtypography
\usepackage{xcolor}         % colors
\usepackage[most]{tcolorbox} % colored boxes
\usepackage{listings}        % code blocks
\usepackage{algorithm}       % algorithms
\usepackage{algpseudocode}   % algorithmic pseudocode

% Listings style for code snippets
\lstset{
  basicstyle=\ttfamily\small,
  columns=fullflexible,
  frame=single,
  breaklines=true,
  showstringspaces=false
}

% Physics Insight box (use whenever you ask to "physics insight" a passage)
\newtcolorbox{physicsinsightbox}{
  enhanced,
  breakable,
  colback=blue!4,
  colframe=blue!55!black,
  boxrule=0.6pt,
  arc=1.2mm,
  left=1.5mm,right=1.5mm,top=1.2mm,bottom=1.2mm,
  fonttitle=\bfseries,
  title={Physics Insight -- Optional}
}

% Mathematical Requirements box (use whenever you ask for "mathematical requirements")
\newtcolorbox{mathrequirementsbox}{
  enhanced,
  breakable,
  colback=red!4,
  colframe=red!65!black,
  boxrule=0.6pt,
  arc=1.2mm,
  left=1.5mm,right=1.5mm,top=1.2mm,bottom=1.2mm,
  fonttitle=\bfseries,
  title={Mathematical Requirements -- Optional}
}

% TL;DR box (use whenever you ask for a short takeaway / summary)
\newtcolorbox{tldrbox}{
  enhanced,
  breakable,
  colback=green!4,
  colframe=green!55!black,
  boxrule=0.6pt,
  arc=1.2mm,
  left=1.5mm,right=1.5mm,top=1.2mm,bottom=1.2mm,
  fonttitle=\bfseries,
  title={TL;DR - IMPORTANT - DO NOT SKIP}
}

% Extra Information box (use for side notes / extra context)
% NOTE: define an explicit purple to avoid class/package color-name clashes.
\definecolor{extrainfopurple}{RGB}{110,55,160}
\newtcolorbox{extrainfobox}{
  enhanced,
  breakable,
  colback=extrainfopurple!5,
  colframe=extrainfopurple!80!black,
  boxrule=0.6pt,
  arc=1.2mm,
  left=1.5mm,right=1.5mm,top=1.2mm,bottom=1.2mm,
  fonttitle=\bfseries,
  title={Extra Information -- Optional}
}

% Examples box
\newtcolorbox{examplebox}{
  enhanced,
  breakable,
  colback=pink!10,
  colframe=pink!60!black,
  boxrule=0.6pt,
  arc=1.2mm,
  left=1.5mm,right=1.5mm,top=1.2mm,bottom=1.2mm,
  fonttitle=\bfseries,
  title={Example}
}

% Computational & Numerical Issues box
\newtcolorbox{compnumissuesbox}{
  enhanced,
  breakable,
  colback=yellow!10,
  colframe=yellow!60!black,
  boxrule=0.6pt,
  arc=1.2mm,
  left=1.5mm,right=1.5mm,top=1.2mm,bottom=1.2mm,
  fonttitle=\bfseries,
  title={Computational \& Numerical Issues}
}

% Note. For the workshop paper template, both \title{} and \workshoptitle{} are required, with the former indicating the paper title shown in the title and the latter indicating the workshop title displayed in the footnote. 
\title{Formatting Instructions For NeurIPS 2026}


% The \author macro works with any number of authors. There are two commands
% used to separate the names and addresses of multiple authors: \And and \AND.
%
% Using \And between authors leaves it to LaTeX to determine where to break the
% lines. Using \AND forces a line break at that point. So, if LaTeX puts 3 of 4
% authors names on the first line, and the last on the second line, try using
% \AND instead of \And before the third author name.


\author{%
  Ignacio E. Adrogué\thanks{Use footnote for providing further information
    about author (webpage, alternative address)---\emph{not} for acknowledging
    funding agencies.} \\
  Department of Physics\\
  ETHZ\\
  Zürich, Switzerland \\
  \texttt{iadrogue@ethz.ch} \\
  % examples of more authors
  % \And
  % Coauthor \\
  % Affiliation \\
  % Address \\
  % \texttt{email} \\
  % \AND
  % Coauthor \\
  % Affiliation \\
  % Address \\
  % \texttt{email} \\
  % \And
  % Coauthor \\
  % Affiliation \\
  % Address \\
  % \texttt{email} \\
  % \And
  % Coauthor \\
  % Affiliation \\
  % Address \\
  % \texttt{email} \\
}


\begin{document}



\maketitle

% Table of contents
\setcounter{tocdepth}{3}
\tableofcontents
\medskip\hrule\medskip

\begin{abstract}
  Explication pas-à-pas des contributions du papier qui devrait permettre même à ceux n'ayant jamais été confronté à du Contrôle Optimal de lire le papier et essayer de comprendre le code, en particulier les examples avec Hamiltonien sans forme fermée. 
\end{abstract}


\section{Part I --- Pontryagin's Maximum Principle}

\subsection{Problem Setup}
\label{sec:problem-setup}

\medskip\hrule\medskip

\subsubsection{The World We Are Describing}

Before writing a single equation, let us be precise about what kind of physical situation we are modeling. We have a \textbf{system} — a rocket, a robot arm, a spacecraft, a bicycle — whose instantaneous condition is completely described by a vector of numbers. We call this vector the \textbf{state} and write it as:

$$z_x(t) \in \mathbb{R}^n$$

The index $n$ is the \textbf{state dimension}. For a particle moving in 3D, $n = 6$ (three positions, three velocities). For a robot arm with five joints, $n = 10$ (five angles, five angular velocities). The state evolves continuously in time over a fixed horizon $t \in [0, T]$.

The state does not evolve on its own — we are allowed to \textbf{intervene}. At each instant, we choose an \textbf{input} to the system, which we call the control:

$$u(t) \in U \subset \mathbb{R}^m$$

The set $U$ is the \textbf{admissible control set} — the physical constraints on what you can actually do. A rocket can only burn fuel in one direction ($U$ is a bounded cone). A steering wheel has a maximum turning angle ($U$ is a closed interval). We require $U$ to be compact (closed and bounded), which is physically natural — no real actuator can produce infinite force.

\medskip\hrule\medskip

\subsubsection{The Dynamics}

The \textbf{law of motion} — the analogue of Newton's second law — takes the form of a controlled ordinary differential equation:

$$\dot{z}_x = f(t, z_x, u), \qquad z_x(0) = x$$

where $f : [0,T] \times \mathbb{R}^n \times U \to \mathbb{R}^n$ is the \textbf{dynamics function}. It tells you, given the current state and the current control input, how fast the state is changing right now.

This is the master equation of the entire framework. Notice what it encodes: the future evolution of $z_x$ depends only on the present state $z_x(t)$ and the present control $u(t)$ — not on anything that came before. This is the \textbf{Markov property}, and it deserves a moment of careful thought because it may seem to conflict with your physical intuition.

\begin{physicsinsightbox}
In Newtonian mechanics, a particle's future trajectory depends on both its current position \textbf{and} its current velocity — you might object that this is non-Markovian, that inertia means the past matters. But this is only an apparent contradiction. The resolution is that \textbf{velocity is part of the state}. When we write $x = (q, \dot{q})$, stacking position and velocity into a single state vector, the evolution $\ddot{q} = F(q, \dot{q})/m$ becomes first-order in this enlarged state — and is then Markovian. The Markov property is not a physical restriction; it is a statement about whether you have chosen a sufficiently complete description of the system. Any finite-order ODE can be written in this first-order Markovian form by enlarging the state space. So when we write $\dot{z}_x = f(t, z_x, u)$, we are implicitly assuming the state $x$ already contains all the information needed to predict the future — positions, velocities, and whatever else is physically relevant.

This structure is exactly the same as Hamilton's equations $\dot{q} = \partial \mathcal{H}/\partial p$, $\dot{p} = -\partial \mathcal{H}/\partial q$ — those are a special case with $x = (q,p)$ and $u$ absent. The Hamiltonian formulation of classical mechanics is already in the form $\dot{x} = f(x)$, and it is Markovian precisely because the canonical momenta $p$ encode the velocity information.
\end{physicsinsightbox}

\begin{mathrequirementsbox}
We impose standard regularity conditions on $f$: it is continuously differentiable in both arguments, and Lipschitz in $z_x$ uniformly over $u$. The Lipschitz condition is the minimal requirement guaranteeing that the ODE has a unique solution for any measurable $u(t)$ — this is the Carathéodory existence theorem. Physically, Lipschitz continuity means the dynamics cannot change infinitely fast in response to a small perturbation in state.
\end{mathrequirementsbox}

\medskip\hrule\medskip

\subsubsection{The Objective: What Does ``Optimal'' Mean?}

Having defined the system, we now need to define what we are trying to achieve. We assign a \textbf{scalar cost} to every possible trajectory and control pair. The cost has two parts.

The first is the \textbf{running cost}:

$$L : [0,T] \times \mathbb{R}^n \times U \to \mathbb{R}$$

$L(t, z_x, u)$ measures the instantaneous penalty at time $t$ — fuel consumption rate, deviation from a desired path, actuator effort. The second is the \textbf{terminal cost}:

$$G : \mathbb{R}^n \to \mathbb{R}$$

$G(z_x(T))$ measures how good or bad the final state is. Together they define the \textbf{objective functional}:

\begin{equation}
\min_{u \in U} \int_0^T L(s, z_x, u)\, ds + G(z_x(T))
\label{eq:ocp-objective}
\end{equation}

\begin{physicsinsightbox}
This is a \textbf{functional} — a map from an entire control trajectory $u(\cdot)$ to a single real number. It is the direct analogue of the \textbf{action functional} in classical mechanics:

$$S[q] = \int_0^T \mathcal{L}(q(t), \dot{q}(t))\, dt$$

The analogy is deep. In classical mechanics you extremize $S$ over paths $q(t)$ with fixed endpoints. In optimal control you minimize $J$ over controls $u(t)$ subject to the dynamics constraint. The variational structure is identical.
\end{physicsinsightbox}

\begin{mathrequirementsbox}
\textbf{A word on differentiating functionals — for the picky mathematician.} The expression $\delta J / \delta u(t)$ can look suspicious. A colleague might argue: "$u(t)$ is a function of $t$, so differentiating $J$ with respect to $u(t)$ is just differentiating with respect to $t$ in disguise." This is not correct, and the distinction is important.

The functional $J[u]$ takes an entire function $u(\cdot)$ as its argument. To differentiate it, we do not differentiate with respect to the variable $t$ — we differentiate with respect to the \textbf{shape of the function} $u$. The precise definition is the \textbf{Fréchet derivative} (or, in the physicist's language, the functional derivative). We perturb the control by a small test function $v(\cdot)$ scaled by $\varepsilon$:

$$J[u + \varepsilon v] = J[u] + \varepsilon \int_0^T \frac{\delta J}{\delta u(t)} v(t)\, dt + O(\varepsilon^2)$$

The coefficient $\delta J / \delta u(t)$ in this expansion is the \textbf{functional derivative} of $J$ with respect to $u$ at time $t$. It is a function of $t$, not a derivative with respect to $t$. The two objects live in entirely different spaces: $d/dt$ acts on functions of time and produces their rate of change; $\delta / \delta u(t)$ acts on functionals of functions and produces how sensitive the functional is to a localized perturbation at time $t$.

% To make this concrete: consider the simple functional $J[u] = \int_0^T u(t)^2\, dt$. Perturbing:

% $$J[u + \varepsilon v] = \int_0^T (u(t) + \varepsilon v(t))^2\, dt = J[u] + 2\varepsilon \int_0^T u(t) v(t)\, dt + O(\varepsilon^2)$$

% Reading off the linear term: $\delta J / \delta u(t) = 2u(t)$. This is simply $2u$ evaluated at time $t$ — it has nothing to do with $du/dt$. The functional derivative is a pointwise operation on the function $u$, not a time derivative.

% This is the mathematical machinery that makes the variational derivation rigorous. When we write $\delta S = 0$ and integrate by parts, we are computing Fréchet derivatives of a constrained functional, not differentiating with respect to $t$.

To make this concrete, consider the functional
\[
J[u] = \int_0^T u(t)^2 \, dt.
\]

We compute its derivative by perturbing $u$ in the direction of an arbitrary test function $v(\cdot)$:
\[
J[u + \varepsilon v]
= \int_0^T (u(t) + \varepsilon v(t))^2 \, dt
= J[u] + 2\varepsilon \int_0^T u(t)v(t)\,dt + \varepsilon^2 \int_0^T v(t)^2\,dt.
\]

Taking $\varepsilon \to 0$, the \emph{Fréchet derivative} of $J$ at $u$ in the direction $v$ is
\[
DJ[u](v) = 2 \int_0^T u(t)v(t)\,dt.
\]

This object is a \emph{linear functional} in $v$, not a function of $t$. In other words, the derivative of $J$ is the mapping
\[
v \mapsto DJ[u](v).
\]

To express this derivative in a pointwise form, we use the standard $L^2$ inner product
\[
\langle a, b \rangle = \int_0^T a(t)b(t)\,dt.
\]
Then the derivative can be written as
\[
DJ[u](v) = \langle 2u, v \rangle.
\]

By definition, the \emph{functional derivative} (or gradient) $\frac{\delta J}{\delta u}(t)$ is the function that satisfies
\[
DJ[u](v) = \int_0^T \frac{\delta J}{\delta u}(t)\, v(t)\, dt
\quad \text{for all } v.
\]

Comparing the two expressions, we identify
\[
\frac{\delta J}{\delta u}(t) = 2u(t).
\]

\paragraph{Remark.}
The distinction is important: the Fréchet derivative $DJ[u]$ is a linear operator, while $\frac{\delta J}{\delta u}(t)$ is its representation as a function under the chosen inner product. This identification depends on the underlying function space (here $L^2$). In informal physics notation, the two are often conflated.
\end{mathrequirementsbox}

\medskip\hrule\medskip

\subsubsection{The Optimal Control Problem --- Formal Statement}

We seek a control function $u^\star : [0,T] \to U$ such that:

$$\min_{u \in U} \int_0^T L(s, z_x, u)\, ds + G(z_x(T))$$

subject to:

\begin{equation}
\dot{z}_x = f(t, z_x, u), \qquad z_x(0) = x, \qquad u(t) \in U \subset \mathbb{R}^m \text{ for all } t
\label{eq:ocp-dynamics}
\end{equation}

The object $u^\star(\cdot)$ is called the \textbf{optimal open-loop control}. The word \textbf{open-loop} is deliberate and important: $u^\star(t)$ is a function of time alone, pre-computed from the initial condition $x$. It is a fixed schedule of inputs — like a pre-programmed rocket burn sequence — that you execute without looking at the actual state of the system as it evolves. If the real system deviates from the expected trajectory due to noise or model error, an open-loop controller does not react.

This stands in contrast to a \textbf{feedback control law} (also called a closed-loop control), which would take the form $u^\star(t, z_x(t))$ — a function of both time and the current state. A feedback law is far more robust, but also far harder to compute. We will see in Part II that the HJB approach naturally produces feedback laws, while PMP naturally produces open-loop trajectories. Bridging the two is a central theme of this document.

The corresponding trajectory $z_x(\cdot)$ obtained by integrating the dynamics under $u^\star$ is called the \textbf{optimal trajectory}. The pair $(z_x, u^\star)$ is called the \textbf{optimal pair}.

\begin{mathrequirementsbox}
We allow $u(t)$ to be any \textbf{measurable} function of time — this is the appropriate function space. We do not require continuity or smoothness of $u$, because in many physical problems the optimal control is discontinuous (ex: bang-bang control, where you switch between maximum and minimum thrust instantaneously).
\end{mathrequirementsbox}

\medskip\hrule\medskip

\begin{mathrequirementsbox}
\subsubsection{The Space of Controls and Why This Is Hard}

The difficulty of this problem, compared to finite-dimensional optimization, is that the decision variable $u(\cdot)$ lives in an \textbf{infinite-dimensional function space}. You are not choosing a vector of $m$ numbers — you are choosing a function on $[0,T]$, which is an uncountably infinite collection of numbers, one for each instant of time, all coupled together through the dynamics ODE.

Standard calculus and gradient descent operate in $\mathbb{R}^m$. To optimize over function spaces, you need the \textbf{calculus of variations} — the machinery of variations, functional derivatives, and adjoint equations. This is what the next subsection develops. The costate $p_x(t)$ will emerge naturally as the functional derivative of $J$ with respect to a perturbation of the trajectory, and the PMP will be the Euler-Lagrange equations of this constrained variational problem.
\end{mathrequirementsbox}

\medskip\hrule\medskip

\begin{mathrequirementsbox}
\subsubsection{Regularity Assumptions --- All in One Place}

To summarize cleanly, we assume throughout Part I: $f$ is $C^1$ in $(t,z_x,u)$ and Lipschitz in $z_x$ uniformly in $u$, ensuring existence and uniqueness of solutions to the dynamics ODE. $L$ and $G$ are $C^1$ in all arguments. $U$ is compact. The optimal pair $(z_x, u^\star)$ exists — we do not prove existence here (that requires Filippov's theorem and convexity of the velocity set, which is a separate story), but we assume it and derive necessary conditions it must satisfy. These necessary conditions are PMP. This is the same epistemological posture as classical mechanics: you assume extremizers exist (Hamilton's principle), and Euler-Lagrange equations tell you what they look like. You don't prove that a classical path exists before deriving Newton's laws.
\end{mathrequirementsbox}


\medskip\hrule\medskip

\subsection{Deriving the equations}
\label{sec:augmented-action}

\medskip\hrule\medskip

\subsubsection{Why We Need an Augmented Functional}

We want to minimize $J[u]$ subject to the constraint $\dot{z}_x = f(t, z_x, u)$ holding at every instant. This is not a finite-dimensional constrained optimization where you have $k$ scalar constraints and introduce $k$ Lagrange multipliers. The constraint $\dot{z}_x(t) = f(t, z_x(t), u(t))$ must hold for every $t \in [0,T]$ — it is an \textbf{infinite-dimensional constraint}, one scalar equation per spatial dimension per instant of time.

The natural generalization of a Lagrange multiplier to this setting is a \textbf{multiplier field}: a time-dependent vector function $p_x(t) \in \mathbb{R}^n$ that enforces the constraint at every $t$ simultaneously. Just as in finite dimensions the Lagrangian $\mathcal{L} = f_0 + \lambda \cdot c$ adds the constraint $c = 0$ weighted by $\lambda$ to the objective $f_0$, here we add the dynamics constraint weighted by $p_x(t)$ and integrated over all time. The result is the \textbf{augmented action}:

\begin{equation}
S[z_x, u, p_x] = \int_0^T \left[ L(t, z_x, u) + \langle p_x(t),\, \dot{z}_x(t) - f(t, z_x, u) \rangle \right] dt + G(z_x(T))
\label{eq:augmented-action}
\end{equation}

Let us read this carefully. The inner product $\langle p_x(t), \dot{z}_x(t) - f(t, z_x, u)\rangle$ is zero whenever the dynamics constraint is satisfied. So on any trajectory that obeys the dynamics, $S$ reduces to exactly $J[u]$. The multiplier field $p_x(t)$ does not change the value of the functional on feasible trajectories — it only penalizes infeasible ones, and its role is to generate the correct optimality conditions when we vary $S$ freely without explicitly enforcing the constraint. This is precisely the Lagrange multiplier philosophy, lifted to function space.

\begin{physicsinsightbox}
Physically, $p_x(t)$ has a beautiful interpretation already visible from its structure: it is the \textbf{generalized momentum conjugate to the state} $z_x$, in the same sense that in classical mechanics the canonical momentum $p = \partial \mathcal{L}/\partial \dot{q}$ is conjugate to $q$. We will see this emerge automatically from the variational computation.
\end{physicsinsightbox}

\medskip\hrule\medskip

\subsubsection{Taking Variations --- The Physics Way}

The following section rederives the PMP equations from first principles. You may rederive them manually, follow the proof or simply trust the process.


\begin{physicsinsightbox}
We now demand that $S$ be stationary under independent variations of all three fields: $z_x \to z_x + \delta z_x$, $u \to u + \delta u$, and $p_x \to p_x + \delta p_x$. The variations $\delta z_x(t)$ and $\delta u(t)$ are arbitrary smooth functions on $[0,T]$, subject only to the boundary condition $\delta z_x(0) = 0$ (we fix the initial condition $z_x(0) = x$ and do not perturb it). The terminal value $\delta z_x(T)$ is \textbf{free} — we impose no constraint on the final state, and the terminal cost $G(z_x(T))$ will generate a boundary condition there.


\paragraph{First-order variation.}
We compute the first variation of the augmented action
\[
S[z_x,u,p_x] = \int_0^T \Big( L(z_x,u) + \langle p_x, \dot z_x - f(z_x,u) \rangle \Big)\,dt + G(z_x(T)).
\]

Let $(\delta z_x, \delta u, \delta p_x)$ be admissible variations. Expanding to first order gives
\[
\delta S
= \int_0^T \Big[
\delta L + \langle \delta p_x, \dot z_x - f \rangle + \langle p_x, \delta \dot z_x - \delta f \rangle
\Big] dt
+ \langle \nabla G(z_x(T)), \delta z_x(T) \rangle.
\]

We linearize each term:
\[
\delta L = \langle \nabla_{z_x} L, \delta z_x \rangle + \langle \nabla_u L, \delta u \rangle,
\]
\[
\delta f = \nabla_{z_x} f \cdot \delta z_x + \nabla_u f \cdot \delta u.
\]

Substituting, we obtain
\[
\delta S
= \int_0^T \Big[
\langle \nabla_{z_x} L, \delta z_x \rangle
+ \langle \nabla_u L, \delta u \rangle
+ \langle \delta p_x, \dot z_x - f \rangle
+ \langle p_x, \delta \dot z_x \rangle
- \langle p_x, \nabla_{z_x} f \cdot \delta z_x \rangle
- \langle p_x, \nabla_u f \cdot \delta u \rangle
\Big] dt
+ \langle \nabla G(z_x(T)), \delta z_x(T) \rangle.
\]

% Expanding $S[z_x + \delta z_x, u + \delta u, p_x + \delta p_x]$ to first order in the variations:

% \begin{align*}
% \delta S
% &= \int_0^T \Bigl[
%     \langle \nabla_{z_x} L,\, \delta z_x \rangle
%   + \langle \nabla_u L,\, \delta u \rangle
%   + \langle \delta p_x,\, \dot{z}_x - f \rangle \\
% &\qquad\qquad\qquad\qquad
%   + \bigl\langle p_x,\, \delta\dot{z}_x - \nabla_{z_x} f \cdot \delta z_x - \nabla_u f \cdot \delta u \bigr\rangle
% \Bigr] \, dt
%   + \langle \nabla G(z_x(T)),\, \delta z_x(T) \rangle .
% \end{align*}

The term $\langle p_x, \delta\dot{z}_x\rangle$ contains a time derivative of the variation. We handle it with \textbf{integration by parts}:

\begin{align*}
\int_0^T \langle p_x(t),\, \delta\dot{z}_x(t) \rangle\, dt
&= \langle p_x(T),\, \delta z_x(T) \rangle
 - \langle p_x(0),\, \delta z_x(0) \rangle
 - \int_0^T \langle \dot{p}_x(t),\, \delta z_x(t) \rangle\, dt .
\end{align*}

Since $\delta z_x(0) = 0$, the boundary term at $t = 0$ vanishes. Substituting back and collecting terms by which variation they multiply:

\begin{align*}
\delta S
&= \int_0^T \langle \delta p_x,\, \dot{z}_x - f \rangle\, dt \\
&\quad + \int_0^T \left\langle \delta z_x,\, \nabla_{z_x} L - \nabla_{z_x} f^\top p_x - \dot{p}_x \right\rangle dt \\
&\quad + \int_0^T \left\langle \delta u,\, \nabla_u L - \nabla_u f^\top p_x \right\rangle dt \\
&\quad + \left\langle \delta z_x(T),\, p_x(T) + \nabla G(z_x(T)) \right\rangle .
\end{align*}

For $\delta S = 0$ to hold for \textbf{all} independent variations $\delta p_x$, $\delta z_x$, $\delta u$, and for the boundary term to vanish, each coefficient must vanish separately. This gives us four conditions, which we now read off one by one.
\end{physicsinsightbox}

\medskip\hrule\medskip

\subsubsection{The PMP equations}

\textbf{The four stationarity conditions, written compactly, are:}
\begin{align*}
\dot{z}_x(t) &= f\bigl(t, z_x(t), u(t)\bigr), & & z_x(0) = x, \\
\dot{p}_x(t) &= \nabla_{z_x} L\bigl(t, z_x(t), u(t)\bigr) - \nabla_{z_x} f\bigl(t, z_x(t), u(t)\bigr)^\top p_x(t), & & p_x(T) = -\nabla G\bigl(z_x(T)\bigr), \\
0 &= \nabla_u L\bigl(t, z_x(t), u(t)\bigr) - \nabla_u f\bigl(t, z_x(t), u(t)\bigr)^\top p_x(t). & &
\end{align*}

\smallskip

These correspond to (i) recovery of the from $\delta p_x$, (ii) the adjoint/costate equation from $\delta z_x$, (iii) stationarity in $u$ from $\delta u$, and (iv) the transversality condition from the boundary term at $t=T$. The minus sign in $p_x(T)$ will shortly be absorbed into the sign convention of the Hamiltonian.

\medskip\hrule\medskip

\begin{tldrbox}
\subsubsection{Recognizing the Hamiltonian}

Now collect the three conditions involving $p_x$ and $z_x$ and observe that they all derive from a single scalar function. Define the \textbf{control Hamiltonian}:
\[
\mathcal{H}(t, z_x, p_x, u) = -\langle p_x,\, f(t, z_x, u)\rangle - L(t, z_x, u).
\label{eq:control-hamiltonian}
\]

Note the sign convention carefully: both terms enter with a minus sign. This is the paper's convention and it is consistent with minimization (rather than maximization) of $J$.

At an optimal control $u^\star$, the PMP relations can be written in compact form:
\begin{equation}
\dot{z}_x(t) = -\nabla_{p_x}\mathcal{H}\bigl(t, z_x(t), p_x(t), u^\star(t)\bigr), \qquad z_x(0)=x,
\label{eq:pmp-state}
\end{equation}
\begin{equation}
\dot{p}_x(t) = \nabla_{z_x}\mathcal{H}\bigl(t, z_x(t), p_x(t), u^\star(t)\bigr), \qquad p_x(T)=\nabla G\bigl(z_x(T)\bigr),
\label{eq:pmp-costate}
\end{equation}
\begin{equation}
0 = \nabla_{u}\mathcal{H}\bigl(t, z_x(t), p_x(t), u^\star(t)\bigr), \qquad \text{(interior optimum)}.
\label{eq:pmp-stationarity}
\end{equation}

The global form of the optimality condition (valid also when $u^\star$ lies on the boundary of $\mathcal{U}$) is
\[
 u^\star(t) \in \arg\max_{u \in \mathcal{U}}\, \mathcal{H}\bigl(t, z_x(t), p_x(t), u\bigr).
\label{eq:pmp-max}
\]

The variational computation has therefore produced, from a single augmented action and a single stationarity demand, all three equations of the paper simultaneously: the state dynamics, the costate dynamics (including transversality), and the optimality condition. The costate $p_x(t)$ emerged not as something postulated but as the natural Lagrange multiplier field enforcing the dynamics constraint — it is the shadow price of the constraint, measuring how much the optimal cost would change if the dynamics were perturbed infinitesimally at time $t$.
\end{tldrbox}


\medskip\hrule\medskip

\subsection{The Equivalent Formulation for Mathematicians}
\label{sec:pmp-theorem}

\medskip\hrule\medskip

\subsubsection{Why Restate What We Already Have?}

\begin{mathrequirementsbox}
The variational derivation above is physically transparent but epistemologically uncomfortable for a mathematician. It assumes smoothness of the optimal control $u^\star$, it treats $\delta u$ as a free variation everywhere on $[0,T]$ without worrying about the constraint $u \in U$, and it derives only a stationarity condition $\nabla_u \mathcal{H} = 0$ rather than a global maximality condition. For controls that hit the boundary of $U$ — which is generically what happens in interesting problems, bang-bang control being the canonical example — $\nabla_u \mathcal{H} = 0$ is simply wrong as a necessary condition. The true necessary condition is the global statement $u^\star \in \arg\max_{u \in U} \mathcal{H}$, and this requires a completely different proof that does not pass through functional derivatives at all.

Pontryagin's Maximum Principle in its full generality is precisely this stronger statement, valid for measurable controls, non-smooth trajectories, and controls on the boundary of $U$. We state it here as a theorem, identify exactly what it gives, and explain what the variational argument captures and what it misses.
\end{mathrequirementsbox}

\medskip\hrule\medskip

\subsubsection{The PMP as a Theorem}

\begin{mathrequirementsbox}
Let $(z_x^\star, u^\star)$ be an optimal pair for the problem:

$$\min_{u \in U} \int_0^T L(s, z_x, u)\, ds + G(z_x(T)), \qquad \dot{z}_x = f(t, z_x, u), \quad z_x(0) = x$$

Under the regularity assumptions of Section~\ref{sec:problem-setup}, there exists an absolutely continuous function $p_x : [0,T] \to \mathbb{R}^n$ such that the following four conditions hold simultaneously almost everywhere on $[0,T]$.

\textbf{Condition 1 — State equation:}
\begin{equation}\tag{3}
\dot{z}_x = -\nabla_{p_x} \mathcal{H}(t, z_x, p_x, u^\star), \qquad z_x(0) = x
\label{eq:pmp-state-tagged}
\end{equation}

\textbf{Condition 2 — Costate equation:}
\begin{equation}\tag{4}
\dot{p}_x = \nabla_{z_x} \mathcal{H}(t, z_x, p_x, u^\star), \qquad p_x(T) = \nabla G(z_x(T))
\label{eq:pmp-costate-tagged}
\end{equation}

\textbf{Condition 3 — Global maximality:}
\begin{equation}\tag{5}
u^\star(t) \in \arg\max_{u \in U}\, \mathcal{H}(t, z_x(t), p_x(t), u) \quad \text{for a.e. } t \in [0,T]
\label{eq:pmp-max-tagged}
\end{equation}

\textbf{Condition 4 — Hamiltonian constancy:}
\begin{equation}
t \mapsto \mathcal{H}(t, z_x(t), p_x(t), u^\star(t)) = \text{const} \quad \text{if } f \text{ and } L \text{ do not depend explicitly on } t
\label{eq:hamiltonian-constancy}
\end{equation}

These are \textbf{necessary conditions} for optimality. They are not sufficient in general. A pair $(z_x^\star, u^\star, p_x)$ satisfying all four conditions is called an \textbf{extremal}, and $p_x$ is the corresponding \textbf{adjoint variable} or \textbf{costate}.
\end{mathrequirementsbox}

\medskip\hrule\medskip

\subsubsection{What the Variational Argument Captured and What It Missed}

\begin{mathrequirementsbox}
The physics-style variational derivation correctly produced the state and costate equations in full generality, and correctly produced $\nabla_u \mathcal{H} = 0$ as a stationarity condition. What it could not produce is the global maximality condition in its full strength.

The gap is this. The variational argument perturbs $u \to u + \varepsilon \delta u$ and demands the first-order term vanish. This is valid only when $u^\star$ lies in the \textbf{interior} of $U$, so that the perturbation $\delta u$ is free in all directions. When $u^\star(t)$ lies on the \textbf{boundary} of $U$ — which is the generic situation in many problems — you cannot perturb freely in all directions, and the condition $\nabla_u \mathcal{H} = 0$ is replaced by the weaker condition that $\mathcal{H}$ is non-increasing as you move away from $u^\star$ into the interior of $U$. The global maximality condition:

$$\mathcal{H}(t, z_x, p_x, u^\star) \geq \mathcal{H}(t, z_x, p_x, u) \quad \forall u \in U$$

subsumes both cases cleanly without needing to distinguish interior from boundary.

The proper proof of PMP uses \textbf{needle variations} (also called Pontryagin variations): instead of perturbing $u$ by a small smooth function everywhere on $[0,T]$, one perturbs it on a single small interval $[t_0 - \varepsilon, t_0]$ by replacing $u^\star$ with an arbitrary fixed value $v \in U$. This localized perturbation is admissible regardless of whether $u^\star(t_0)$ is interior or boundary, and the resulting first-order expansion of $J$ yields the global maximality condition directly. The costate $p_x$ emerges in this proof as the solution to the adjoint equation that makes the first-order variation of $J$ expressible as a single inner product — it is the Riesz representer of the functional derivative of $J$ with respect to needle perturbations.

We do not reproduce the full needle variation proof here — it is lengthy and primarily functional-analytic in character. What matters for our purposes is the output: the state equation, the costate equation, and the global maximality condition, now valid for measurable $u^\star$ and controls on the boundary of $U$.
\end{mathrequirementsbox}

\medskip\hrule\medskip

\subsubsection{The Concavity Condition and the Implication in \eqref{eq:pmp-max-tagged}}

\begin{mathrequirementsbox}
The paper writes \eqref{eq:pmp-max-tagged} as:

$$u^\star \in \arg\max_u\, \mathcal{H}(t, z_x, p_x, u) \implies \nabla_u \mathcal{H}(t, z_x, p_x, u^\star) = 0$$

The implication arrow deserves scrutiny — as you correctly flagged earlier. This implication holds if and only if two conditions are met:

\textbf{First}, $u^\star$ must lie in the interior of $U$. If $u^\star$ is on the boundary, the gradient $\nabla_u \mathcal{H}$ at $u^\star$ need not be zero — it only needs to point outward or be zero.

\textbf{Second}, $\mathcal{H}$ must be differentiable in $u$ at $u^\star$. This is guaranteed by our $C^1$ assumption on $f$ and $L$.

When both conditions hold, the global maximum of $\mathcal{H}$ over $U$ at an interior point satisfies the first-order condition $\nabla_u \mathcal{H} = 0$ by elementary calculus. But this is only a necessary condition for the maximum — it finds all critical points, not just maxima. To guarantee that the solution to $\nabla_u \mathcal{H} = 0$ is indeed the \textbf{global maximum} and not a saddle point or minimum, we need:

$$\nabla^2_{uu} \mathcal{H}(t, z_x, p_x, u^\star) \prec 0$$

that is, the Hessian of $\mathcal{H}$ with respect to $u$ must be \textbf{negative definite} at $u^\star$ — equivalently, $\mathcal{H}$ must be \textbf{strictly concave} in $u$. Under strict concavity, $\nabla_u \mathcal{H} = 0$ is both necessary and sufficient for a global maximum, and the maximizer is unique. This is the condition that makes the implicit solve of Part IV well-posed, and it is why the paper's implicit Hamiltonian approach requires it as a standing assumption.
\end{mathrequirementsbox}

\medskip\hrule\medskip

\subsubsection{PMP as a Two-Point Boundary Value Problem}

\begin{physicsinsightbox}
Conditions (3) and (4) together form a system of $2n$ first-order ODEs in the $2n$-dimensional space of $(z_x, p_x)$:

$$\frac{d}{dt}\begin{pmatrix} z_x \\ p_x \end{pmatrix} = \begin{pmatrix} -\nabla_{p_x}\mathcal{H} \\ \nabla_{z_x}\mathcal{H} \end{pmatrix}$$

This is a \textbf{Hamiltonian system} — the same mathematical structure as Hamilton's equations in classical mechanics, with $z_x$ playing the role of generalized coordinates and $p_x$ playing the role of generalized momenta. The function $\mathcal{H}$ is the Hamiltonian, and the flow it generates in the $(z_x, p_x)$ phase space is a symplectic flow — it preserves the symplectic form $dp_x \wedge dz_x$.

The boundary conditions however are \textbf{split}: $z_x(0) = x$ is known at the left endpoint, while $p_x(T) = \nabla G(z_x(T))$ is specified at the right endpoint. This is a \textbf{two-point boundary value problem} (BVP), not an initial value problem. You cannot simply integrate forward from $t=0$ because you do not know $p_x(0)$ — you only know $p_x(T)$, and even that depends on $z_x(T)$ which you don't know until after integration.

This is the structural reason PMP is computationally hard. The standard approach — the \textbf{shooting method} — guesses $p_x(0)$, integrates the full system (3)–(4) forward, checks whether $p_x(T) = \nabla G(z_x(T))$ is satisfied, and iterates. The sensitivity of this shooting problem is governed by the \textbf{monodromy matrix} of the Hamiltonian flow, whose eigenvalues grow exponentially with $T$ in chaotic or unstable systems. Small errors in the guess for $p_x(0)$ amplify exponentially, making the shooting method unreliable for long horizons or high-dimensional systems. This is the precise failure mode of PMP that the HJB approach in Part II is designed to circumvent — at the cost of the curse of dimensionality.
\end{physicsinsightbox}


\medskip\hrule\medskip

\subsection{The Explicit Hamiltonian Case and When PMP Can Be Solved}
\label{sec:explicit-case}

\medskip\hrule\medskip

\subsubsection{What ``Explicit'' Means and Why It Matters}

\begin{tldrbox}
Having derived the PMP conditions \eqref{eq:pmp-state-tagged}--\eqref{eq:pmp-max-tagged}, the question of practical solvability reduces to a single bottleneck: can you evaluate the maximality condition \eqref{eq:pmp-max-tagged} in closed form? That is, given the current state $z_x$, costate $p_x$, and time $t$, can you write down an explicit formula:

$$u^\star = g(t, z_x, p_x)$$

for the optimal control as a function of the known quantities? If yes, the Hamiltonian is \textbf{explicit} — you can substitute $u^\star$ back into the PMP state and costate equations and reduce the PMP system to a closed ODE in $(z_x, p_x)$ alone, which you can then attempt to solve numerically via a shooting method. If no — if $\arg\max_{u \in U} \mathcal{H}$ has no closed form — the Hamiltonian is \textbf{implicit}, and this substitution cannot be made. The entire PMP-based computational pipeline breaks at this step.

This distinction is the central axis around which the paper's contribution turns. We develop the explicit case fully here, both to build intuition and to make precise exactly what fails in the implicit case.
\end{tldrbox}

\medskip\hrule\medskip

\subsubsection{The Canonical Explicit Case: Control-Affine Systems with Quadratic Cost}

\begin{examplebox}
The most important and widely studied class of explicit problems has two ingredients working together. First, the dynamics are \textbf{control-affine} — the vector field $f$ is linear in $u$:

$$f(t, z_x, u) = a(t, z_x) + B(t, z_x)\, u$$

where $a : [0,T] \times \mathbb{R}^n \to \mathbb{R}^n$ is the drift and $B : [0,T] \times \mathbb{R}^n \to \mathbb{R}^{n \times m}$ is the control matrix. Second, the running cost is \textbf{quadratic in $u$}:

$$L(t, z_x, u) = \frac{1}{2} u^\top R\, u + C(t, z_x)$$

where $R \in \mathbb{R}^{m \times m}$ is symmetric positive definite (so $R \succ 0$) and $C$ is a state-dependent term that does not involve $u$. The positive definiteness of $R$ is the precise condition that makes $L$ strictly convex in $u$, which will translate into strict concavity of $\mathcal{H}$ in $u$.

The control Hamiltonian is:

\begin{align*}
\mathcal{H}(t, z_x, p_x, u)
&= -\langle p_x,\, a + Bu \rangle - \frac{1}{2} u^\top R u - C(t, z_x) \\
&= -\langle p_x, a \rangle - \langle p_x, Bu \rangle - \frac{1}{2} u^\top R u - C(t, z_x) \\
&= -\langle p_x, a \rangle - \langle B^\top p_x, u \rangle - \frac{1}{2} u^\top R u - C(t, z_x) .
\end{align*}

This is a \textbf{quadratic function of $u$}, with leading term $-\frac{1}{2} u^\top R u$. Since $R \succ 0$, the leading term is strictly negative definite in $u$, so $\mathcal{H}$ is strictly concave in $u$ globally — not just locally. The unique global maximum over $\mathbb{R}^m$ is found by setting the gradient to zero:

$$\nabla_u \mathcal{H} = -B^\top p_x - Ru = 0$$

Solving immediately:

$$u^\star = -R^{-1} B(t, z_x)^\top p_x$$

This is the explicit formula $u^\star = g(t, z_x, p_x)$. The optimal control is a \textbf{linear function of the costate} $p_x$, with the matrix $R^{-1} B^\top$ acting as the gain. If additionally $U = \mathbb{R}^m$ (unconstrained control), this is the full answer. If $U$ is a compact convex set, the constrained maximizer is the projection of $-R^{-1}B^\top p_x$ onto $U$, which is also explicit when $U$ is a simple set like a box or ball.
\end{examplebox}

\medskip\hrule\medskip

\subsubsection{Substituting Back: The Reduced Hamiltonian System}

\begin{examplebox}
Now substitute $u^\star = -R^{-1}B^\top p_x$ back into \eqref{eq:pmp-state-tagged} and \eqref{eq:pmp-costate-tagged}. The state equation \eqref{eq:pmp-state-tagged} becomes:

$$\dot{z}_x = a(t, z_x) - B(t, z_x) R^{-1} B(t, z_x)^\top p_x$$

The costate equation \eqref{eq:pmp-costate-tagged} becomes:

$$\dot{p}_x = \nabla_{z_x} \mathcal{H}\big|_{u = u^\star} = -\nabla_{z_x} a^\top p_x + \nabla_{z_x}[B^\top p_x] R^{-1} B^\top p_x - \nabla_{z_x} C + \frac{1}{2}\nabla_{z_x}[p_x^\top B R^{-1} B^\top p_x]$$

In the special case where $a$ and $B$ are both independent of $z_x$ — the \textbf{linear time-varying} system $\dot{z}_x = A(t) z_x + B(t) u$ with $a = Az_x$ — this simplifies dramatically and the system (3)–(4) becomes a linear BVP, solvable via the \textbf{Riccati equation}. This is the foundation of Linear Quadratic Regulator (LQR) theory, the most celebrated explicit optimal control solution.

In the nonlinear case, the substitution still closes the system into a $2n$-dimensional autonomous ODE in $(z_x, p_x)$, which can be attacked numerically via the shooting method with a well-defined right-hand side at every step.
\end{examplebox}

\medskip\hrule\medskip

\subsubsection{The Bang-Bang Case: Explicit but Non-Smooth}

\begin{examplebox}
A second important explicit case arises when $\mathcal{H}$ is \textbf{linear} in $u$ and $U$ is a hypercube $U = \{u : |u_i| \leq u_{\max}\}$. Write:

$$\mathcal{H}(t, z_x, p_x, u) = \langle \sigma(t, z_x, p_x),\, u \rangle + (\text{terms independent of } u)$$

where $\sigma(t, z_x, p_x) = -\nabla_u(\langle p_x, f\rangle + L)$ is the \textbf{switching function}. Since $\mathcal{H}$ is linear in $u$, its maximum over the hypercube is achieved at a vertex — a corner where each component of $u$ is at its extreme value. The maximizer is:

$$u^\star_i(t) = u_{\max} \cdot \text{sign}(\sigma_i(t, z_x, p_x))$$

This is \textbf{bang-bang control} — the optimal control switches discontinuously between its extreme values, with the switching times determined by the zeros of $\sigma$. It is explicit: given $(z_x, p_x)$, you can compute $u^\star$ immediately from the sign of $\sigma$. No equation needs to be solved. But $u^\star$ is discontinuous, which is why the PMP needs to be stated for measurable controls rather than smooth ones — this is exactly the case the needle variation proof covers and the variational argument misses.
\end{examplebox}

\medskip\hrule\medskip

\subsubsection{What Explicit Buys You Computationally}

\begin{tldrbox}
In both the quadratic and bang-bang cases, the key computational consequence is the same: the PMP system (3)–(5) reduces to a \textbf{closed system of ODEs} in $(z_x, p_x)$ with a well-defined right-hand side that can be evaluated pointwise. Specifically, at any $(t, z_x, p_x)$ you can compute $u^\star$ in $O(m)$ or $O(m^2)$ operations (matrix-vector multiply or sign evaluation), then evaluate $f$ and $\nabla_{z_x}\mathcal{H}$ to advance the ODE integrator one step. The entire shooting method requires only ODE integration plus a nonlinear root-finding iteration on $p_x(0)$.

This is what previous neural network approaches to value function approximation inherited as a requirement. When they parameterize $\phi_\theta(t, x)$ and need to generate training trajectories, they must integrate the dynamics forward under $u^\star$, which requires evaluating $u^\star = g(t, z_x, \nabla_{z_x}\phi_\theta)$ at every integration step. If $g$ has no closed form, this evaluation is impossible — or rather, it requires an inner iterative solve at every step of every trajectory during training, which is precisely the implicit Hamiltonian problem the paper addresses.
\end{tldrbox}

\medskip\hrule\medskip

\subsubsection{The Failure Mode: When No Closed Form Exists}

\begin{tldrbox}
Consider now a dynamics function $f(t, z_x, u)$ that is nonlinear in $u$ in a way that cannot be inverted analytically, or a running cost $L(t, z_x, u)$ involving $u$ through a transcendental function. The condition $\nabla_u \mathcal{H} = 0$ then becomes a nonlinear equation in $u$ with no algebraic solution. For example, if:

$$f(t, z_x, u) = a(t, z_x) + b(t, z_x) \sin(u)$$

then $\nabla_u \mathcal{H} = 0$ gives:

$$-\langle p_x,\, b(t, z_x) \rangle \cos(u) - \nabla_u L = 0$$

which is a transcendental equation in $u$ with no closed-form solution in general. Or consider aerodynamic models where the lift and drag forces depend on the angle of attack $u$ through tabulated experimental data or high-fidelity simulation — no formula for $u^\star$ exists at all.

In these cases, $\mathcal{H}$ is \textbf{implicit}: the Hamiltonian is well-defined and evaluable, and $u^\star$ is well-defined by the maximality condition (guaranteed to exist and be unique if $\mathcal{H}$ is strictly concave in $u$), but $u^\star$ cannot be written as an explicit function of $(t, z_x, p_x)$. It can only be characterized as the solution to:

$$F(u;\, t, z_x, p_x) := \nabla_u \mathcal{H}(t, z_x, p_x, u) = 0$$

This is a fixed-point or root-finding problem \textbf{that must be solved numerically — at every instant, at every training step, inside every trajectory simulation}. Managing this inner solve efficiently and differentiating through it to train $\phi_\theta$ is exactly the technical core of the paper, developed in Part IV.
\end{tldrbox}

\medskip\hrule\medskip

\subsection{PMP as a Boundary Value Problem and Its Failure Mode}
\label{sec:pmp-bvp-failure}

\medskip\hrule\medskip

\subsubsection{Recasting PMP as a Shooting Problem}

As discussed above, PMP produces a two-point boundary value problem. Let us now make this completely concrete and show precisely where and why it breaks down computationally, both analytically and through a code implementation.

\begin{mathrequirementsbox}
In the explicit case, substituting $u^\star = g(t, z_x, p_x)$ into the PMP state and costate equations closes the system. Define the augmented phase space vector $\xi = (z_x, p_x) \in \mathbb{R}^{2n}$ and the closed right-hand side:

$$F(t, \xi) = \begin{pmatrix} -\nabla_{p_x}\mathcal{H}(t, z_x, p_x, g(t, z_x, p_x)) \\ \nabla_{z_x}\mathcal{H}(t, z_x, p_x, g(t, z_x, p_x)) \end{pmatrix}$$

The PMP system is then simply $\dot{\xi} = F(t, \xi)$, a standard ODE. The boundary conditions are:

$$z_x(0) = x \qquad \text{(known, left endpoint)}$$
$$p_x(T) = \nabla G(z_x(T)) \qquad \text{(known only implicitly, right endpoint)}$$

You know $n$ components of $\xi$ at $t=0$ and the other $n$ components only at $t=T$ through a condition that itself depends on the unknown $z_x(T)$. To integrate the ODE forward from $t=0$, you need all $2n$ components of $\xi(0)$ — but $p_x(0)$ is unknown.

The \textbf{shooting method} converts this BVP into a root-finding problem. Introduce the shooting variable $q \in \mathbb{R}^n$ as a guess for the initial costate $p_x(0) = q$. For any $q$, integrate the ODE forward:

$$\dot{\xi}(t) = F(t, \xi(t)), \qquad \xi(0) = (x, q)$$

to obtain $\xi(T; q) = (z_x(T; q), p_x(T; q))$. Define the \textbf{shooting residual}:

$$R(q) = p_x(T; q) - \nabla G(z_x(T; q))$$
\end{mathrequirementsbox}

The true $p_x(0)$ is the root of $R(q) = 0$. This is an $n$-dimensional nonlinear root-finding problem, typically solved by Newton's method:

$$q^{(k+1)} = q^{(k)} - \left[\frac{\partial R}{\partial q}(q^{(k)})\right]^{-1} R(q^{(k)})$$

\begin{tldrbox}
The Jacobian $\partial R / \partial q$ is the \textbf{monodromy matrix} — it measures how the endpoint $(z_x(T), p_x(T))$ varies when the initial costate $p_x(0)$ is perturbed. Computing it requires integrating $2n$ additional variational ODEs alongside the main system, one for each column of the Jacobian. The total cost of one Newton step is therefore $O(2n \times 2n)$ ODE integrations — already $O(n^2)$ in state dimension, before even accounting for the nonlinear iterations needed for convergence.
\end{tldrbox}

\medskip\hrule\medskip

\subsubsection{The Exponential Sensitivity: A Precise Statement}

\begin{compnumissuesbox}
The fundamental difficulty is not just computational cost but numerical conditioning. The monodromy matrix $M(T) = \partial \xi(T) / \partial \xi(0)$ satisfies the variational equation:

$$\dot{M}(t) = \nabla_\xi F(t, \xi(t))\, M(t), \qquad M(0) = I_{2n}$$

For a Hamiltonian system, $M(t)$ is a \textbf{symplectic matrix} at all times — its eigenvalues come in pairs $(\lambda, 1/\lambda)$. If the Hamiltonian flow has any instability (which is generic for nonlinear systems and long horizons), some eigenvalues of $M(T)$ will have magnitude $e^{\lambda_{\max} T}$ where $\lambda_{\max} > 0$ is the largest Lyapunov exponent of the flow. The condition number of the shooting Jacobian therefore grows as:

$$\kappa\left(\frac{\partial R}{\partial q}\right) \sim e^{\lambda_{\max} T}$$

exponentially in the time horizon $T$. This means that a perturbation of size $\varepsilon$ in the guess $q$ produces an error of size $\varepsilon \cdot e^{\lambda_{\max} T}$ in the terminal condition — the shooting problem becomes exponentially ill-conditioned as $T$ grows. In double precision arithmetic with $\varepsilon \approx 10^{-16}$, the shooting method loses all accuracy once $\lambda_{\max} T \gtrsim 36$.

This is not a deficiency of the shooting method specifically — it reflects the genuine difficulty of the BVP, which is that the stable and unstable manifolds of the Hamiltonian flow become exponentially separated over long times, and finding the trajectory that connects a given initial state to a given terminal costate condition requires exponential precision in the initial guess.
\end{compnumissuesbox}

\medskip\hrule\medskip

\subsubsection{Dependence on Initial Conditions}

\begin{tldrbox}
The second failure mode is more subtle and directly relevant to the paper's motivation. The shooting solution $q^\star(x)$ — the correct initial costate as a function of the initial state $x$ — is a complicated nonlinear function that must be recomputed from scratch for every new initial condition $x$. There is no reuse: solving the shooting problem for $x_1$ gives you no information about the solution for $x_2 \neq x_1$.

This is the open-loop nature of PMP made precise. The output of the shooting method is a single optimal trajectory $(z_x^\star(t), u^\star(t))$ valid only for the specific initial condition $x$ it was solved for. If you want a \textbf{feedback law} — a map $u^\star(t, x)$ valid for all initial conditions in some region of state space — you would need to solve the shooting problem for a dense grid of initial conditions, which costs $O(N^n)$ shooting solves for $N$ points per dimension. This is a different manifestation of the curse of dimensionality from the one afflicting HJB, but it is equally crippling for high-dimensional systems.

The value function approach of Part II sidesteps this entirely: $\phi(t, x)$ encodes the optimal cost for all initial conditions simultaneously, and $\nabla_x \phi(t, x)$ gives the optimal costate — and hence the optimal control — at any state instantaneously, without re-solving any BVP.
\end{tldrbox}

\medskip\hrule\medskip

\subsubsection{Code Implementation: Where Shooting Breaks}

We implement the shooting method for a simple but illustrative optimal control problem: the nonlinear pendulum with energy cost. The state is $z_x = (\theta, \dot{\theta}) \in \mathbb{R}^2$, the control is torque $u \in \mathbb{R}$, the dynamics are $\dot{z}_x = (\dot{\theta},\, -\sin\theta + u)$, the running cost is $L = \frac{1}{2}(z_x^\top Q z_x + r u^2)$, and the terminal cost is $G(z_x(T)) = \frac{1}{2} z_x(T)^\top Q_f z_x(T)$. This is an explicit Hamiltonian problem — $u^\star = -p_{x,2}/r$ — so the shooting method is applicable in principle. We use it to demonstrate the sensitivity failure as $T$ and $n$ grow.

% \begin{lstlisting}[language=Python]
% import numpy as np
% from scipy.integrate import solve_ivp
% from scipy.optimize import fsolve
% import matplotlib.pyplot as plt
% import time

% # -- Problem parameters -------------------------------------------------------
% r   = 1.0                          # control weight
% Q   = np.eye(2)                    # state cost
% Qf  = 5.0 * np.eye(2)             # terminal cost

% def dynamics(t, z, u):
%     return np.array([z[1], -np.sin(z[0]) + u])

% def optimal_u(p):
%     # explicit formula: u* = -p[1] / r  (from grad_u H = 0)
%     return -p[1] / r

% def hamiltonian_rhs(t, xi):
%     """
%     Full 2n-dimensional Hamiltonian system RHS.
%     xi = [z_x (n), p_x (n)]
%     """
%     n  = len(xi) // 2
%     z  = xi[:n]
%     p  = xi[n:]
%     u  = optimal_u(p)

%     # state equation:  dz/dt = f(z, u)
%     dz = dynamics(t, z, u)

%     # costate equation: dp/dt = grad_z H
%     # H = -<p, f> - L,  L = 0.5*(z'Qz + r*u^2)
%     # grad_z H = -grad_z f^T p - grad_z L
%     # grad_z f = [[0,1],[-cos(z[0]),0]]
%     grad_z_f = np.array([[0.0, 1.0], [-np.cos(z[0]), 0.0]])
%     grad_z_L = Q @ z
%     dp = -grad_z_f.T @ p - grad_z_L  # note: sign consistent with eq (4)

%     return np.concatenate([dz, dp])

% def shooting_residual(q, x0, T):
%     """
%     Given initial costate guess q = p_x(0), integrate forward and
%     return the transversality residual R(q) = p_x(T) - grad G(z_x(T)).
%     """
%     xi0  = np.concatenate([x0, q])
%     sol  = solve_ivp(hamiltonian_rhs, [0, T], xi0,
%                      method='RK45', rtol=1e-10, atol=1e-12, dense_output=False)
%     n    = len(x0)
%     zT   = sol.y[:n,  -1]
%     pT   = sol.y[n:,  -1]
%     return pT - Qf @ zT   # transversality: p_x(T) = grad G(z_x(T)) = Qf @ z_x(T)

% # -- Experiment 1: sensitivity to initial costate guess ----------------------
% x0    = np.array([1.0, 0.0])   # initial state: pendulum displaced, at rest
% T     = 3.0
% q_true, info, ier, msg = fsolve(
%     shooting_residual, np.zeros(2), args=(x0, T),
%     full_output=True, xtol=1e-12
% )
% print(f"Shooting converged: {ier==1}, residual norm: {np.linalg.norm(info['fvec']):.2e}")

% # Perturb initial costate by epsilon and measure terminal error
% epsilons   = np.logspace(-10, -1, 50)
% term_errors = []
% for eps in epsilons:
%     q_perturbed  = q_true + eps * np.array([1.0, 1.0]) / np.sqrt(2)
%     res          = shooting_residual(q_perturbed, x0, T)
%     term_errors.append(np.linalg.norm(res))

% plt.figure(figsize=(7, 4))
% plt.loglog(epsilons, term_errors, 'b-', linewidth=2)
% plt.loglog(epsilons, epsilons,    'k--', label='slope 1 (linear)')
% plt.xlabel('Perturbation size eps in p_x(0)')
% plt.ylabel('Terminal residual ||R(q* + eps*delta)||')
% plt.title('Shooting sensitivity: pendulum, T=3')
% plt.legend(); plt.grid(True); plt.tight_layout()
% plt.savefig('shooting_sensitivity.png', dpi=150)
% plt.close()

% # -- Experiment 2: condition number growth with horizon T --------------------
% def monodromy_condition_number(x0, T, q0=None):
%     """
%     Estimate condition number of shooting Jacobian dR/dq via finite differences.
%     """
%     n   = len(x0)
%     if q0 is None:
%         q0 = np.zeros(n)
%     J   = np.zeros((n, n))
%     R0  = shooting_residual(q0, x0, T)
%     h   = 1e-6
%     for i in range(n):
%         dq      = np.zeros(n); dq[i] = h
%         R1      = shooting_residual(q0 + dq, x0, T)
%         J[:, i] = (R1 - R0) / h
%     return np.linalg.cond(J)

% horizons = np.linspace(0.5, 6.0, 20)
% cond_numbers = []
% for T_val in horizons:
%     try:
%         q0, _, ier, _ = fsolve(
%             shooting_residual, np.zeros(2), args=(x0, T_val),
%             full_output=True, xtol=1e-10
%         )
%         if ier == 1:
%             kappa = monodromy_condition_number(x0, T_val, q0)
%             cond_numbers.append((T_val, kappa))
%     except Exception:
%         pass

% Tvals, kappas = zip(*cond_numbers)
% plt.figure(figsize=(7, 4))
% plt.semilogy(Tvals, kappas, 'r-o', linewidth=2, markersize=5)
% plt.xlabel('Time horizon T')
% plt.ylabel('Condition number kappa(dR/dq)')
% plt.title('Monodromy matrix conditioning vs. horizon')
% plt.grid(True); plt.tight_layout()
% plt.savefig('monodromy_conditioning.png', dpi=150)
% plt.close()

% # -- Experiment 3: cost of solving for each new initial condition ------------
% initial_conditions = [np.array([np.cos(th), np.sin(th)])
%                       for th in np.linspace(0, np.pi, 20)]
% T_fixed = 2.0
% solve_times = []
% for x_init in initial_conditions:
%     t0 = time.time()
%     fsolve(shooting_residual, np.zeros(2), args=(x_init, T_fixed), xtol=1e-10)
%     solve_times.append(time.time() - t0)

% print(f"\nShooting solve times for 20 initial conditions (T={T_fixed}):")
% print(f"  Mean: {np.mean(solve_times)*1000:.1f} ms")
% print(f"  Total: {np.sum(solve_times)*1000:.1f} ms")
% print(f"  Each solve is independent -- zero reuse across initial conditions.")
% \end{lstlisting}

\medskip\hrule\medskip

\subsubsection{What the Code Demonstrates}

The three experiments expose the three independent failure modes of PMP as a computational tool.

\textbf{Experiment 1} shows the sensitivity to the initial costate guess. The terminal residual grows linearly with the perturbation size — as expected, since the monodromy matrix is the linearized sensitivity — but the proportionality constant grows with $T$, so for long horizons even machine-precision perturbations in $q$ produce $O(1)$ errors in the terminal condition. In practice this means Newton's method for the shooting problem fails to converge unless initialized extremely close to the true $p_x(0)$, which you do not know without having already solved the problem.

\textbf{Experiment 2} shows the condition number of the shooting Jacobian growing exponentially with $T$. Plotting $\kappa(\partial R / \partial q)$ on a log scale against $T$ reveals a straight line — exponential growth — consistent with the Lyapunov exponent analysis above. For the nonlinear pendulum, the instability of trajectories near the separatrix drives this growth. For higher-dimensional systems with larger Lyapunov exponents, the conditioning deteriorates even faster.

\textbf{Experiment 3} exposes the open-loop nature of PMP directly: each initial condition requires a completely independent shooting solve, with no information shared between them. Twenty initial conditions require twenty full shooting iterations. A feedback law over a grid of $N^n$ initial conditions would require $N^n$ such solves — the curse of dimensionality in PMP's clothing.

Together, these three experiments make precise the computational argument for moving beyond PMP to the value function approach of Part II: exponential sensitivity to initialization, exponential conditioning with horizon, and no reuse across initial conditions. The HJB equation, to which we now turn, solves all three problems in principle — at the cost of a different exponential, the curse of dimensionality in state space dimension.

\pagebreak






% \medskip\hrule\medskip

\section{Part II --- The Hamilton-Jacobi-Bellman Equation}

\subsection{Bellman's Principle of Optimality and the Value Function}
\label{sec:hjb}

\medskip\hrule\medskip

\subsubsection{A Different Question Entirely}

\begin{tldrbox}
The PMP approach of Part I asks: given a specific initial condition $x$, what is the optimal control trajectory $u^\star(t)$? It answers by solving a boundary value problem in the phase space $(z_x, p_x)$, producing a single optimal trajectory valid only for that $x$. As shown in Section~\ref{sec:pmp-bvp-failure}, this is computationally fragile and produces no reusable structure across initial conditions.


Part II asks a fundamentally different question: what is the optimal cost achievable from \textbf{every possible initial condition simultaneously}? The answer to this question is a single scalar function defined over the entire state space, and once you have it, the optimal control at any state at any time is immediately recoverable by differentiation. This function is the \textbf{value function}, and it is the central object of dynamic programming.
\end{tldrbox}

\begin{physicsinsightbox}
The shift in perspective is profound. Instead of thinking about optimal trajectories — one-dimensional curves in state space — we think about an optimal \textbf{field} over state space. This is analogous to the shift in classical mechanics from thinking about individual particle trajectories to thinking about Hamilton's principal function $S(q, t)$, which encodes all trajectories simultaneously as its level surfaces. The analogy is not coincidental — they are mathematically the same object.
\end{physicsinsightbox}

\medskip\hrule\medskip

\subsubsection{The Value Function --- Definition}

Recall the optimal control problem. For an initial condition $x \in \mathbb{R}^n$ at time $t = 0$, the optimal cost is:

$$\min_{u \in U} \int_0^T L(s, z_x, u)\, ds + G(z_x(T))$$

subject to $\dot{z}_x = f(s, z_x, u)$, $z_x(0) = x$. Now generalize: instead of always starting at $t = 0$, consider starting from state $y$ at an arbitrary intermediate time $t \in [0, T]$. Define the \textbf{value function}:

\begin{equation}
\phi(t, x) = \inf_{u \in U} \left\{ \int_t^T L(s, z_x, u)\, ds + G(z_x(T)) \right\}
\label{eq:value-function}
\end{equation}

subject to $\dot{z}_x = f(s, z_x, u)$, $z_x(t) = x$.

This is the paper's notation: $\phi$ for the value function, with arguments $(t, x)$ where $t$ is the current time and $x$ is the current state. The infimum is over all admissible controls on the interval $[t, T]$. The subscript on $z_x$ now carries a double meaning: $z_x$ is the state trajectory starting from $x$ at time $t$, so more precisely $z_x(s; t, x)$, but we suppress the dependence on $(t,x)$ for notational cleanliness.

Several immediate observations. First, $\phi(t, x)$ is a function of \textbf{two arguments}: time $t$ and state $x$. It lives on the $(n+1)$-dimensional space $[0,T] \times \mathbb{R}^n$.

\begin{tldrbox}
Second, the terminal condition is fixed: $\phi(T, x) = G(x)$ for all $x$, since starting at the terminal time leaves no control freedom and the cost is just the terminal penalty. Third, the value function at $t=0$ evaluated at the specific initial condition $x_0$ recovers the original optimal cost: $\phi(0, x_0) = \min_u J[u]$. The value function is a global object. It encodes the answer to the optimal control problem for every possible initial condition simultaneously. This is its power and its computational burden.
\end{tldrbox}

\medskip\hrule\medskip

\subsubsection{Bellman's Principle of Optimality}

The value function satisfies a fundamental \textbf{self-consistency condition} known as Bellman's principle of optimality. It can be stated informally as:

\begin{tldrbox}
\emph{An optimal policy has the property that, whatever the initial state and initial decision, the remaining decisions must constitute an optimal policy with regard to the state resulting from the first decision.}
\end{tldrbox}

More precisely, consider the optimal trajectory starting from $(t, x)$. Pick any intermediate time $t + \delta t$ with $\delta t > 0$ small. The trajectory from $t$ to $T$ is optimal. Therefore the portion from $t + \delta t$ to $T$ — starting from the state $z_x(t + \delta t)$ reached by the optimal control on $[t, t+\delta t]$ — must itself be optimal for the subproblem starting at $(t + \delta t, z_x(t + \delta t))$. If it were not, you could replace it with a better policy and reduce the total cost, contradicting the optimality of the original trajectory.

This self-consistency is not a theorem to be proved from outside — it is a consequence of the additive separable structure of the cost functional $J$. Because cost accumulates as an integral, the cost from $t$ to $T$ decomposes as cost from $t$ to $t+\delta t$ plus cost from $t+\delta t$ to $T$. This additive decomposition is what makes dynamic programming work. It would fail for non-additive costs such as the maximum of $L$ over the trajectory, or a ratio of integrals.

Formally, Bellman's principle states:

\begin{equation}
\phi(t, x) = \inf_{u \in U} \left\{ \int_t^{t+\delta t} L(s, z_x, u)\, ds + \phi(t + \delta t,\, z_x(t + \delta t)) \right\}
\label{eq:bellman-principle}
\end{equation}

for any $\delta t > 0$. The value function at $(t, x)$ equals the minimum over controls on the short interval $[t, t+\delta t]$ of the running cost on that interval plus the value function at the state reached at the end of the interval. The future optimal cost from $t + \delta t$ onward is exactly $\phi(t+\delta t, z_x(t+\delta t))$ — this is Bellman's principle packaging the entire future optimization into the single scalar $\phi$ evaluated at the next state.

\medskip\hrule\medskip

\subsubsection{From Bellman's Principle to the HJB Equation}
This subsection derives the HJB equation from first principles, which is not strictly necessary to follow the paper, but is quite instructive and easy to read.

\begin{mathrequirementsbox}
We now take $\delta t \to 0$ to convert Bellman's principle into a PDE. Assume $\phi$ is $C^1$ in both arguments — we will comment on when this fails shortly.

Expand $\phi(t + \delta t, z_x(t + \delta t))$ to first order in $\delta t$. Using the chain rule:

\begin{align*}
\phi(t + \delta t, z_x(t + \delta t))
&= \phi(t, x) + \delta t\, \frac{\partial \phi}{\partial t}(t, x) + \delta t \left\langle \nabla_x \phi(t, x),\, \dot{z}_x(t) \right\rangle + O(\delta t^2) \\
&= \phi(t, x) + \delta t\, \frac{\partial \phi}{\partial t}(t, x) + \delta t \left\langle \nabla_x \phi(t, x),\, f(t, x, u) \right\rangle + O(\delta t^2) .
\end{align*}

Also expand the running cost integral over the short interval:

$$\int_t^{t+\delta t} L(s, z_x, u)\, ds = \delta t\, L(t, x, u) + O(\delta t^2)$$

Substituting into Bellman's principle:

\begin{align*}
\phi(t, x)
&= \inf_{u \in U} \Bigl\{ \delta t\, L(t, x, u) + \phi(t, x) + \delta t\, \frac{\partial \phi}{\partial t} + \delta t \left\langle \nabla_x \phi,\, f(t, x, u) \right\rangle + O(\delta t^2) \Bigr\} .
\end{align*}

Subtract $\phi(t, x)$ from both sides — it does not depend on $u$ and passes through the infimum:

$$0 = \inf_{u \in U} \left\{ \delta t\, L(t, x, u) + \delta t\, \frac{\partial \phi}{\partial t} + \delta t \left\langle \nabla_x \phi,\, f(t, x, u) \right\rangle \right\} + O(\delta t^2)$$

Divide by $\delta t$ and take $\delta t \to 0$:

$$0 = \frac{\partial \phi}{\partial t}(t, x) + \inf_{u \in U} \left\{ L(t, x, u) + \left\langle \nabla_x \phi(t, x),\, f(t, x, u) \right\rangle \right\}$$

\begin{tldrbox}
    
Rearranging:

$$-\frac{\partial \phi}{\partial t}(t, x) = \inf_{u \in U} \left\{ L(t, x, u) + \left\langle \nabla_x \phi(t, x),\, f(t, x, u) \right\rangle \right\}$$

with terminal condition $\phi(T, x) = G(x)$. This is the \textbf{Hamilton-Jacobi-Bellman equation}. It is a first-order PDE in $(n+1)$ variables, solved \textbf{backward in time} from the terminal condition $\phi(T, \cdot) = G$.
\end{tldrbox}
\end{mathrequirementsbox}

\medskip\hrule\medskip

\subsubsection{The Hamiltonian Appears Again}

\begin{tldrbox}
Define the \textbf{value-function Hamiltonian} — the same object that will appear in the paper's notation — by performing the optimization over $u$ inside the HJB equation:

\begin{equation}
H(t, x, p) = \sup_{u \in U} \left[ -\langle p, f(t, x, u) \rangle - L(t, x, u) \right] = \sup_{u \in U} \mathcal{H}(t, x, p, u)
\label{eq:optimized-hamiltonian}
\end{equation}

This is the \textbf{Legendre-Fenchel transform} of $L$ with respect to $u$, composed with $f$ (see Section~\ref{sec:legendre-fenchel-viewpoint}). With this definition, the HJB equation takes the compact form:

\begin{equation}
-\frac{\partial \phi}{\partial t}(t, x) + H(t, x, \nabla_x \phi(t, x)) = 0, \qquad \phi(T, x) = G(x)
\label{eq:hjb-hamiltonian-form}
\end{equation}

\begin{physicsinsightbox}
This is the paper's equation relating $\phi$ to $H$. It has the identical structure to the \textbf{Hamilton-Jacobi equation} of classical mechanics:

$$\frac{\partial S}{\partial t} + H\!\left(q, \frac{\partial S}{\partial q}\right) = 0$$

where $S(q,t)$ is Hamilton's principal function. The value function $\phi$ is the optimal control analogue of Hamilton's principal function — both satisfy the same PDE, with the Hamiltonian playing the same role in both. This is not a coincidence: classical mechanics is a special case of optimal control where there are no external controls and the cost is the action.
\end{physicsinsightbox}

The optimal feedback control is then recovered as:

\begin{equation}
u^\star(t, x) = \arg\sup_{u \in U}\, \mathcal{H}(t, x, \nabla_x \phi(t, x), u)
\label{eq:hjb-feedback}
\end{equation}

This is a \textbf{true feedback law}: given any state $x$ at any time $t$, compute $\nabla_x \phi(t, x)$ and maximize $\mathcal{H}$ over $u$. No BVP, no shooting, no dependence on initial conditions. The entire optimal policy is encoded in the single scalar field $\phi(t, x)$.
\end{tldrbox}

\medskip\hrule\medskip




\subsection{The Curse of Dimensionality and the Failure of Grid-Based HJB}
\label{sec:curse}

\medskip\hrule\medskip

\subsubsection{Why the HJB PDE Cannot Be Solved Directly in High Dimensions}

The HJB equation, first written in the compact Hamiltonian form in \eqref{eq:hjb-hamiltonian-form},

\begin{equation}
    -\frac{\partial \phi}{\partial t}(t, x) + H(t, x, \nabla_x \phi(t, x)) = 0, \qquad \phi(T, x) = G(x),
\end{equation}

\begin{tldrbox}
is a first-order PDE defined over the $(n+1)$-dimensional domain $[0,T] \times \mathbb{R}^n$. To solve it numerically in the classical way, one discretizes this domain with a grid of $N$ points per spatial dimension, producing a grid of total size $N^n$ in the state space alone. For $n = 2$ and $N = 100$, this is $10^4$ points --- entirely manageable. For $n = 6$ and $N = 100$, this is $10^{12}$ points --- already beyond the memory of any existing computer. For $n = 20$, the number $N^{20}$ is astronomically large regardless of $N$. This exponential growth of computational cost with state dimension is the \textbf{curse of dimensionality}, a term coined by Bellman himself when he recognized that dynamic programming, despite its conceptual elegance, was computationally intractable for high-dimensional systems.

\begin{compnumissuesbox}
It is important to be precise about what \textit{kind} of exponential this is, and how it differs from the exponential that afflicts PMP. The PMP shooting problem becomes exponentially ill-conditioned as the \textit{time horizon} $T$ grows, through the Lyapunov exponent mechanism described in Section~\ref{sec:pmp-bvp-failure}. The HJB grid discretization becomes exponentially costly as the \textit{state dimension} $n$ grows, through the volume of the grid. These are independent exponentials in different variables. A problem with short horizon and high state dimension would suffer from HJB's curse but not PMP's; a problem with long horizon and low state dimension would suffer from PMP's conditioning but not HJB's grid cost. The regime the paper targets --- high-dimensional state space \textit{and} long horizon --- is where \textit{both} classical approaches fail simultaneously. This is the precise gap the paper fills.
\end{compnumissuesbox}
\end{tldrbox}

\medskip\hrule\medskip

\subsubsection{The Cost Decomposition: Memory, Computation, and Time}

\begin{compnumissuesbox}
To make the curse concrete, consider three separate costs that each scale as $N^n$.

\textbf{Memory cost.} Storing the value function $\phi(t, x)$ on the grid requires one floating-point number per grid point per time step. With $N^n$ spatial points and $N_t$ time steps, total storage is $N_t \cdot N^n$ numbers. At 8 bytes per double-precision float, storing a grid for $n = 10$, $N = 10$, $N_t = 100$ requires $100 \times 10^{10} \times 8 \approx 8 \times 10^{12}$ bytes --- 8 terabytes.

\textbf{Computation cost per time step.} At each time step, the HJB equation must be evaluated at every grid point. Computing $\nabla_x \phi$ at a grid point via finite differences requires $2n$ neighboring evaluations. The $\inf_u$ optimization at each point requires an inner loop over $u$. Total operations per time step: $O(N^n \cdot n)$. Total over all time steps: $O(N_t \cdot N^n \cdot n)$.

\textbf{Interpolation cost.} Recovering the optimal control at a state $x$ that does not lie exactly on a grid point requires $n$-linear interpolation, which costs $O(2^n)$ operations --- itself exponential in $n$.

All three costs share the same fundamental structure: they grow exponentially with $n$ regardless of any other problem parameter.
\end{compnumissuesbox}

\medskip\hrule\medskip

\subsubsection{Code Implementation: Timing the Explosion}

We implement a grid-based HJB solver for the same nonlinear pendulum problem used in Section~\ref{sec:pmp-bvp-failure}, then extend to higher state dimensions to measure the explosion directly.
Code not shown for brevity
% \begin{verbatim}
% import numpy as np
% import time
% import matplotlib.pyplot as plt
% from itertools import product

% # -- Problem parameters (same as the PMP failure-mode section) ----------------
% T    = 2.0
% r    = 1.0
% Q    = None          # will be eye(n) for each n
% Qf   = None          # will be 5*eye(n)

% def build_grid(n, N, x_min=-3.0, x_max=3.0):
%     """
%     Build an n-dimensional grid with N points per dimension.
%     Returns the grid as a list of 1D arrays (one per dimension)
%     and the total number of points.
%     """
%     axes       = [np.linspace(x_min, x_max, N) for _ in range(n)]
%     total_pts  = N ** n
%     return axes, total_pts

% def hjb_grid_cost(n, N, N_t=50):
%     """
%     Measure memory and timing cost of setting up and doing one
%     backward time step of HJB on an n-dimensional grid with N pts/dim.
%     Returns (memory_bytes, time_seconds) or None if too large.
%     """
%     total_pts = N ** n
%     memory_bytes = total_pts * 8   # float64

%     # refuse if > 4 GB
%     if memory_bytes > 4e9:
%         return memory_bytes, None

%     t0 = time.time()

%     # allocate value function grid
%     shape = tuple([N] * n)
%     phi   = np.zeros(shape, dtype=np.float64)

%     # terminal condition: phi(T, x) = G(x) = 0.5 * x'Qf x
%     # for efficiency, build using meshgrid only up to n=5
%     if n <= 5:
%         axes = [np.linspace(-3.0, 3.0, N) for _ in range(n)]
%         grids = np.meshgrid(*axes, indexing='ij')
%         phi_T = sum(5.0 * 0.5 * g**2 for g in grids)
%         phi[:] = phi_T
%     else:
%         # too large to meshgrid: just fill with zeros as proxy
%         phi[:] = 0.0

%     # one backward Euler step of HJB (simplified: just evaluate
%     # the finite-difference gradient at every grid point)
%     dt     = T / N_t
%     dx     = 6.0 / (N - 1)

%     # compute grad_x phi via central differences at every interior point
%     # only dimension 0 for timing purposes (representative cost)
%     slc_fwd = [slice(None)] * n
%     slc_bwd = [slice(None)] * n
%     slc_fwd[0] = slice(2, N)
%     slc_bwd[0] = slice(0, N-2)
%     grad0 = (phi[tuple(slc_fwd)] - phi[tuple(slc_bwd)]) / (2 * dx)

%     elapsed = time.time() - t0
%     return memory_bytes, elapsed

% # -- Experiment: sweep over (n, N) -------------------------------------------
% dimensions = [2, 3, 4, 5, 6, 8, 10]
% N_pts      = [20, 15, 10, 8, 6, 4, 3]   # decreasing N as n grows

% print(f"{'n':>4} {'N':>4} {'Grid points':>15}"
%       f" {'Memory (MB)':>12} {'Time (s)':>10}")
% print("-" * 55)

% results = []
% for n, N in zip(dimensions, N_pts):
%     total_pts    = N ** n
%     mem_bytes, t = hjb_grid_cost(n, N)
%     mem_mb       = mem_bytes / 1e6
%     t_str        = f"{t:.4f}" if t is not None else "OOM"
%     print(f"{n:>4} {N:>4} {total_pts:>15,} {mem_mb:>12.1f} {t_str:>10}")
%     results.append((n, N, total_pts, mem_mb, t))

% # -- Plot: grid points vs state dimension -------------------------------------
% ns    = [r[0] for r in results]
% pts   = [r[2] for r in results]
% times = [r[4] for r in results if r[4] is not None]
% ns_t  = [r[0] for r in results if r[4] is not None]

% fig, axes = plt.subplots(1, 2, figsize=(12, 4))

% axes[0].semilogy(ns, pts, 'b-o', linewidth=2, markersize=6)
% axes[0].set_xlabel('State dimension $n$')
% axes[0].set_ylabel('Total grid points $N^n$')
% axes[0].set_title('Grid size explosion with dimension')
% axes[0].grid(True)

% axes[1].semilogy(ns_t, times, 'r-o', linewidth=2, markersize=6)
% axes[1].set_xlabel('State dimension $n$')
% axes[1].set_ylabel('Wall time (s), one HJB step')
% axes[1].set_title('Computation time vs.\ dimension')
% axes[1].grid(True)

% plt.tight_layout()
% plt.savefig('hjb_cod.png', dpi=150)
% plt.close()

% print("\nConclusion: both grid size and compute time grow as N^n.")
% print("For n=10, N=10: grid has 10^10 points, requiring ~80 GB.")
% print("Classical HJB is computationally dead beyond n ~ 6.")
% \end{verbatim}

\medskip\hrule\medskip

\subsubsection{What the Code Shows}

\begin{compnumissuesbox}
Running the above produces a table and two plots that make the curse of dimensionality visible as a physical phenomenon rather than an abstract asymptotic statement.

For $n = 2$ with $N = 20$, the grid has $400$ points and the computation takes milliseconds --- this is the textbook regime where HJB is tractable and widely used in robotics and control engineering. For $n = 6$ with $N = 6$, the grid already has $46{,}656$ points and the coarseness of the grid ($N=6$ points per dimension) means the value function is barely resolved --- yet even this crude discretization begins to strain memory. For $n = 10$ with $N = 3$, the grid has $3^{10} = 59{,}049$ points but each dimension has only 3 points, making the approximation essentially useless. A meaningful discretization with $N = 20$ for $n = 10$ would require $20^{10} \approx 10^{13}$ points, or roughly 80 terabytes of memory.

The timing plot shows the computation time growing exponentially on a log scale --- a straight line on the semilog axes --- confirming the theoretical $O(N^n)$ scaling. The slope of the line is $\log N$ per unit of $n$, consistent with the grid volume formula.
\end{compnumissuesbox}

\medskip\hrule\medskip

\subsection{Final insights on HJB}
\label{sec:HJB-final-insights}

\subsubsection{The Situation So Far: Two Failures, One Gap}

At this point in the document we have established the following precise picture.

\textbf{PMP} converts the optimal control problem into a $2n$-dimensional Hamiltonian BVP. It produces open-loop trajectories for individual initial conditions. Its computational cost is independent of the grid but grows exponentially with horizon $T$ through the conditioning of the monodromy matrix. It requires a closed-form expression for $u^\star = g(t, z_x, p_x)$ to close the ODE system, which fails for implicit Hamiltonians.

\textbf{HJB} converts the optimal control problem into an $(n+1)$-dimensional PDE. It produces global feedback laws valid for all initial conditions. Its computational cost is independent of horizon but grows as $N^n$ with state dimension $n$. It also requires a closed-form expression for the $\sup_u \mathcal{H}$ to evaluate the right-hand side of the PDE at each grid point, which again fails for implicit Hamiltonians.

\textbf{Neither approach dominates the other.} They trade one exponential for another in different variables. For the high-dimensional, long-horizon problems the paper targets, both fail simultaneously and for different reasons.

The next part, after a math interlude for the curious reader, develops the PMP--HJB bridge and shows how this identity can be combined with implicit solves and neural parameterization. 



\subsubsection{Open-Loop vs.\ Feedback: A Precise Comparison}

\begin{tldrbox}
We can now state precisely the relationship between the PMP open-loop solution
and the HJB feedback solution for the same problem.

Given an initial condition $x$ at $t=0$, the PMP produces a trajectory
$(z_x^\star(t), p_x^\star(t))$ and an open-loop control $u^\star(t) =
g(t, z_x^\star(t), p_x^\star(t))$. This control is optimal \textit{only along
the trajectory $z_x^\star$}. If at any time $s \in (0,T)$ the state is perturbed
to $z_x^\star(s) + \varepsilon$, the pre-computed control $u^\star(t)$ for
$t > s$ is no longer optimal for the perturbed state --- it was designed for the
unperturbed trajectory.

The HJB feedback law $u^\star(t,x) = \arg\sup_u \mathcal{H}(t, x, \nabla_x
\phi(t,x), u)$ is optimal for \textit{any} state $x$ at \textit{any} time $t$.
If the system is perturbed at time $s$ to state $\tilde{x}$, the feedback law
immediately provides the optimal action $u^\star(s, \tilde{x})$ for the new
state, without any re-computation. This is the \textbf{closed-loop} property
and it is what makes feedback control robust in practice.

The two are consistent along the optimal trajectory: the open-loop control
$u^\star(t)$ produced by PMP agrees with the feedback law evaluated along
$z_x^\star(t)$:

\begin{equation}
    u^\star_{\text{PMP}}(t) = \arg\sup_{u \in U}\,
    \mathcal{H}(t, z_x^\star(t), \nabla_x\phi(t, z_x^\star(t)), u) = u^\star_{\text{HJB}}(t, z_x^\star(t))
\end{equation}

The precise PMP--HJB identity that underpins this agreement is proved in Part III.
\end{tldrbox}

\subsection{Mathematical subtleties of HJB for those interested}
\label{sec:optimized-hamiltonian}

\medskip\hrule\medskip

\subsubsection{Viscosity Solutions: When $\phi$ Is Not Smooth}

\begin{mathrequirementsbox}
A critical mathematical point must be addressed. The derivation above assumed $\phi \in C^1$. In general, the value function is \textbf{not} smooth — it can develop kinks and discontinuities in its gradient even when the data $f$, $L$, $G$ are smooth. This happens when optimal trajectories cross (forming caustics, as in geometrical optics) or when bang-bang optimal controls create switching surfaces.

At points where $\phi$ is not differentiable, the HJB equation does not hold in the classical sense. The correct notion of solution is the \textbf{viscosity solution}, introduced by Crandall and Lions in 1983. A viscosity solution replaces pointwise derivatives with one-sided test function conditions and is the unique weak solution of the HJB equation under mild assumptions. For our purposes, the key fact is that the value function $\phi$ is always a viscosity solution of the HJB equation, even when not classically differentiable. When $\phi$ is smooth at a point $(t,x)$, the viscosity solution coincides with the classical one. Since neural network approximations of $\phi$ produce smooth functions by construction, they implicitly work in the regime where $\phi$ is smooth — a standing assumption that is reasonable for many practical problems and is implicitly made throughout the paper.
\end{mathrequirementsbox}


\subsubsection{The Legendre--Fenchel Viewpoint}
\label{sec:legendre-fenchel-viewpoint}

\begin{mathrequirementsbox}
The function $H(t, x, p)$ defined in \eqref{eq:H_def} is not just a notational
convenience --- it has deep structure inherited from convex analysis. To see this,
fix $(t, x)$ and think of $H$ as a function of $p$ alone:

\begin{equation}
    H(t, x, p) = \sup_{u \in U} \left[ -\langle p, f(t,x,u) \rangle
    - L(t,x,u) \right].
    \label{eq:H_def}
\end{equation}

This is a \textbf{supremum of affine functions} of $p$ (each $u$ contributes the
affine function $p \mapsto -\langle p, f(t,x,u)\rangle - L(t,x,u)$), which means
$H$ is \textbf{convex in $p$} for any fixed $(t,x)$, regardless of the structure
of $f$ and $L$. This is a fundamental and useful property: the optimized
Hamiltonian is always convex in the costate variable, even for nonlinear and
non-convex problems.

When $f$ is linear in $u$ and $L$ is convex in $u$, equation \eqref{eq:H_def}
is precisely the \textbf{Legendre--Fenchel transform} (convex conjugate) of $L$
with respect to $u$, evaluated at $-B^\top p$ where $B$ is the control matrix.
In the quadratic case $L = \frac{1}{2}u^\top R u + C(t,x)$ with $f = a + Bu$,
the supremum in \eqref{eq:H_def} is achieved at $u^\star = -R^{-1}B^\top p$
(from the explicit-case subsection), giving:

\begin{equation}
    H(t, x, p) = -\langle p, a \rangle - C(t,x)
    + \frac{1}{2} p^\top B R^{-1} B^\top p.
\end{equation}

This is a quadratic function of $p$, as expected from the Legendre transform of
a quadratic. The positive semi-definite term $\frac{1}{2}p^\top B R^{-1} B^\top
p$ is what makes $H$ convex in $p$.

In classical mechanics, the Legendre transform converts the Lagrangian
$\mathcal{L}(q, \dot{q})$ into the Hamiltonian $\mathcal{H}(q, p)$ via $\mathcal{H}
= p \cdot \dot{q} - \mathcal{L}$, where $p = \partial \mathcal{L}/\partial \dot{q}$.
The construction \eqref{eq:H_def} is the direct generalization: $p$ plays the role
of the canonical momentum, $f(t,x,u)$ plays the role of $\dot{q}$, and the
supremum over $u$ plays the role of the Legendre transform. The analogy is
exact when the control $u$ is unconstrained and $L$ is strictly convex in $u$.
\end{mathrequirementsbox}

\medskip\hrule\medskip

\subsubsection{The Verification Theorem}

\begin{mathrequirementsbox}
The statement that $u^\star$ defined by \eqref{eq:feedback_law} is indeed optimal
--- not just a candidate --- is the content of the \textbf{verification theorem}.
It is worth stating precisely because it clarifies the logical relationship
between $\phi$, the HJB equation, and optimality.

\medskip

\noindent\textbf{Theorem (Verification).} \label{thm:verification} \textit{Let $\psi \in C^1([0,T] \times
\mathbb{R}^n)$ be a classical solution to the HJB equation:}
\begin{equation}
    -\frac{\partial \psi}{\partial t}(t,x)
    + H(t, x, \nabla_x \psi(t,x)) = 0,
    \qquad \psi(T,x) = G(x).
\end{equation}
\textit{Then $\psi(t,x) \leq \phi(t,x)$ for all $(t,x)$, with equality
$\psi = \phi$ if there exists a measurable $u^\star(t,x)$ achieving the
$\sup$ in $H$ for all $(t,x)$. In that case, the feedback law
\eqref{eq:feedback_law} is optimal and $\psi$ is the value function.}

\medskip

The proof uses It\^{o}'s lemma (or the chain rule in the deterministic case)
to differentiate $\psi(t, z_x(t))$ along any trajectory, integrate from $t$
to $T$, and use the HJB inequality to bound the result. The key step is:

\begin{equation}
    \frac{d}{ds}\psi(s, z_x(s))
    = \frac{\partial \psi}{\partial t}
    + \langle \nabla_x \psi,\, f(s, z_x, u) \rangle
    = -H(s, z_x, \nabla_x\psi)
    + \langle \nabla_x\psi, f \rangle
    \geq L(s, z_x, u)
\end{equation}

where the last inequality uses the definition of $H$ as a supremum. Integrating
from $t$ to $T$ and using the terminal condition $\psi(T, z_x(T)) = G(z_x(T))$:

\begin{equation}
    G(z_x(T)) - \psi(t,x)
    = \int_t^T \frac{d}{ds}\psi(s, z_x(s))\, ds
    \geq \int_t^T L(s, z_x, u)\, ds
\end{equation}

rearranging:

\begin{equation}
    \psi(t,x) \leq \int_t^T L(s, z_x, u)\, ds + G(z_x(T))
\end{equation}

for any admissible $u$. Taking the infimum over $u$ on the right: $\psi(t,x)
\leq \phi(t,x)$. Equality holds when $u$ achieves the supremum in $H$ at every
instant, which is exactly the feedback law \eqref{eq:feedback_law}. This
confirms that any smooth solution to the HJB equation that admits a measurable
maximizer is the value function, and the maximizer is the optimal feedback
control.

The verification theorem is the rigorous justification for the neural network
approach of Part IV: if we train $\phi_\theta$ to satisfy the HJB equation
approximately, and the approximation error is small, then the feedback law
derived from $\phi_\theta$ via \eqref{eq:feedback_law} is approximately optimal.
The quality of the control policy is directly tied to the quality of the HJB
residual.
\end{mathrequirementsbox}


\pagebreak

\section{Part III --- The PMP--HJB Bridge}

\subsection{The PMP--HJB Bridge via the Method of Characteristics}
\label{sec:bridge}

\medskip\hrule\medskip

\subsubsection{Two Frameworks, One Object}

Parts I and II developed PMP and HJB as independent frameworks, each with its
own derivation, its own computational profile, and its own failure mode. We now
establish rigorously that they are not merely analogous --- they are two
representations of the same underlying mathematical object. The bridge is the
\textbf{method of characteristics} applied to the HJB PDE, which generates
exactly the PMP Hamiltonian system as its characteristic equations. The costate
$p_x$ and the gradient $\nabla_x \phi$ are the same object viewed from the two
frameworks.

This connection is classical --- it was known to Carathéodory, Bellman, and
Pontryagin's school in the 1950s and 1960s. We present it here not as a new
result but as the essential mathematical foundation for the paper's method:
without this identity, one cannot substitute $\nabla_x \phi_\theta$ for $p_x$
in the implicit solve of Part IV.

\medskip\hrule\medskip

\subsubsection{Characteristics of a First-Order PDE}

Recall that the HJB equation is a first-order PDE in $(n+1)$ variables:

\begin{equation}
    \phi_t + H(t, x, \nabla_x \phi) = 0, \qquad \phi(T, x) = G(x).
    \label{eq:hjb_char}
\end{equation}

For a general first-order PDE $\mathcal{F}(t, x, \phi, \phi_t, p) = 0$ where
$p = \nabla_x \phi$, the \textbf{method of characteristics} seeks curves
$(t(s), x(s))$ in the $(t,x)$-space along which the PDE reduces to a system
of ODEs. These curves are the \textbf{characteristics}, and they carry
information about the solution $\phi$ from the terminal condition backward
in time.

For the HJB equation \eqref{eq:hjb_char}, we have $\mathcal{F} = \phi_t +
H(t, x, p) = 0$. The characteristic equations are obtained by differentiating
$\mathcal{F} = 0$ along a curve and demanding consistency. Specifically, define
the characteristic strip $(t(s), x(s), \phi(s), p(s))$ where $p(s) =
\nabla_x \phi(t(s), x(s))$, and demand that both the PDE and its gradient with
respect to $x$ are satisfied along the curve. This yields the
\textbf{characteristic system}:

\begin{align}
    \dot{t}(s)   &= 1
    \label{eq:char_t}\\
    \dot{x}(s)   &= \nabla_p H(t, x, p)
    \label{eq:char_x}\\
    \dot{p}(s)   &= -\nabla_x H(t, x, p)
    \label{eq:char_p}\\
    \dot{\phi}(s) &= \phi_t + \langle \nabla_x\phi,\, \dot{x} \rangle
                  = -H + \langle p,\, \nabla_p H \rangle
    \label{eq:char_phi}
\end{align}

where dots denote derivatives with respect to the curve parameter $s$, and
we use the PDE \eqref{eq:hjb_char} to substitute $\phi_t = -H$ in
\eqref{eq:char_phi}. Since $\dot{t} = 1$ by \eqref{eq:char_t}, the parameter
$s$ is identified with $t$ directly, and we can write everything as ODEs in
$t$.

\medskip\hrule\medskip

\subsubsection{The Characteristics Are the PMP Trajectories}
\begin{mathrequirementsbox}
\noindent\textbf{The characteristics are the PMP trajectories.}

We now identify the characteristic equations \eqref{eq:char_x}--\eqref{eq:char_p}
with the PMP system. Recall from Section~\ref{sec:optimized-hamiltonian} that:

\begin{equation}
    \nabla_p H(t, x, p) = \nabla_p \mathcal{H}(t, x, p, u^\star)
    = -f(t, x, u^\star)
\end{equation}

\begin{equation}
    \nabla_x H(t, x, p) = \nabla_x \mathcal{H}(t, x, p, u^\star)
\end{equation}

where $u^\star = \arg\sup_{u \in U} \mathcal{H}(t, x, p, u)$ and we use the
envelope theorem: since $u^\star$ maximizes $\mathcal{H}$, the derivative of
$H = \sup_u \mathcal{H}$ with respect to a parameter equals the derivative of
$\mathcal{H}$ evaluated at the maximizer, with no contribution from the
variation of $u^\star$ itself. Substituting:

\begin{equation}
    \dot{x} = \nabla_p H = -f(t, x, u^\star)
    \qquad \Longrightarrow \qquad
    \dot{x} = -\nabla_{p_x}\mathcal{H}(t, x, p, u^\star)
    \label{eq:char_pmp3}
\end{equation}

\begin{equation}
    \dot{p} = -\nabla_x H = -\nabla_x \mathcal{H}(t, x, p, u^\star)
    \qquad \Longrightarrow \qquad
    \dot{p} = \nabla_{z_x}\mathcal{H}(t, x, p, u^\star)
    \label{eq:char_pmp4}
\end{equation}

where in \eqref{eq:char_pmp4} we use $-\nabla_x\mathcal{H} =
\nabla_{z_x}\mathcal{H}$ up to the sign convention of the paper's $\mathcal{H}
= -\langle p, f\rangle - L$. These are \textbf{exactly \eqref{eq:pmp-state-tagged} and \eqref{eq:pmp-costate-tagged}
of the paper}: the PMP state equation and costate equation. The characteristics
of the HJB PDE are the optimal trajectories of PMP, running in the
$(x, p)$-phase space under the Hamiltonian $H$.

This is the geometric content of the PMP--HJB equivalence: the HJB PDE
propagates the value function backward in time along curves that are precisely
the forward-time optimal trajectories of PMP. The two frameworks solve the same
problem by propagating information in opposite temporal directions along the
same set of curves.
\end{mathrequirementsbox}

\medskip\hrule\medskip

\subsubsection{Deriving the Identity $p_x = \nabla_x \phi$}

\begin{mathrequirementsbox}
Let $(z_x^\star(t), u^\star(t))$ be an optimal trajectory-control pair, and let
$p_x^\star(t)$ be the associated PMP costate with terminal condition
$p_x^\star(T) = \nabla G(z_x^\star(T))$. Assume $\phi$ is smooth on a
neighborhood of this trajectory.

\medskip

\noindent\textbf{Proposition.} \label{prop:gradient-identity} \textit{For all } $t \in [0,T]$,
\begin{equation}
    p_x^\star(t) = \nabla_x \phi(t, z_x^\star(t)).
    \label{eq:bridge}
\end{equation}

\medskip

\noindent\textbf{Proof.}
\begin{enumerate}
    \item Define two candidate costates along the same optimal state trajectory:
    \[
    q_1(t) := p_x^\star(t), \qquad
    q_2(t) := \nabla_x \phi(t, z_x^\star(t)).
    \]

    \item From the characteristic system of the HJB PDE, $q_2$ satisfies
    \[
    \dot{q}_2(t) = -\nabla_x H\!\left(t, z_x^\star(t), q_2(t)\right).
    \]

    \item From PMP and the envelope relation between $H$ and $\mathcal{H}$,
    $q_1$ satisfies the same ODE:
    \[
    \dot{q}_1(t) = -\nabla_x H\!\left(t, z_x^\star(t), q_1(t)\right).
    \]

    \item At terminal time, both satisfy the same boundary value:
    \[
    q_1(T) = p_x^\star(T) = \nabla G(z_x^\star(T)),
    \]
    while from $\phi(T,x)=G(x)$,
    \[
    q_2(T)=\nabla_x\phi(T,z_x^\star(T))=\nabla G(z_x^\star(T)).
    \]

    \item Hence $q_1$ and $q_2$ solve the same ODE with the same terminal
    condition. By uniqueness of ODE solutions (Picard--Lindel\"of),
    $q_1(t)=q_2(t)$ on $[0,T]$, proving \eqref{eq:bridge}. \hfill $\square$
\end{enumerate}
\end{mathrequirementsbox}

\medskip\hrule\medskip

\subsubsection{What This Identity Gives Us and What It Does Not}

The identity \eqref{eq:bridge} is the mathematical engine of the paper's method,
but it must be understood precisely in terms of what it enables and what it
does not.

\textbf{What it gives.} Along any optimal trajectory starting from $x$,
the PMP costate is the gradient of the value function. This means that if we
have a good approximation $\phi_\theta \approx \phi$, we can substitute
$p_x \approx \nabla_x \phi_\theta(t, z_x)$ and obtain an approximate costate
without solving any BVP. The optimal control is then:

\begin{equation}
    u^\star \approx \arg\sup_{u \in U}\,
    \mathcal{H}(t, z_x, \nabla_x\phi_\theta(t, z_x), u)
\end{equation}

evaluated via the fixed-point iteration when no closed form exists. This is the paper's
method in one equation.

\textbf{What it does not give.} The identity holds \textit{along optimal
trajectories}, and it requires $\phi$ to be the true value function, not an
approximation. During training, $\phi_\theta$ is only an approximation, and the
trajectories generated under $\nabla_x\phi_\theta$ are only approximately
optimal. The training process must therefore be understood as an iterative
procedure: as $\phi_\theta$ improves, the generated trajectories become more
nearly optimal, and the identity \eqref{eq:bridge} is satisfied more accurately,
which in turn improves the training signal. Convergence of this loop ---
and the role of JFB in making it computationally tractable --- is the subject
in the subsequent sections.

\textbf{Why prior work still needed explicit Hamiltonians.} Previous methods
that used the identity \eqref{eq:bridge} still required a closed-form expression
$u^\star = g(t, z_x, \nabla_x\phi_\theta)$ to generate training trajectories.
Substituting $p_x = \nabla_x\phi_\theta$ into the explicit formula $g$ is
straightforward. Without a closed form, the substitution produces an implicit
equation for $u^\star$ that those methods had no machinery to solve inside the
training loop. The paper's contribution is precisely the addition of the implicit
solve via the fixed-point iteration and the JFB gradient approximation that makes
differentiating through that solve tractable.

\medskip\hrule\medskip

\begin{physicsinsightbox}
\subsubsection{The Symplectic Structure of the Characteristic Flow}
\label{subsec:lagrangian-submanifold}

We close this step with a remark on the geometric structure that underlies the
PMP--HJB equivalence, which is important for understanding why the bridge works
so cleanly.

The characteristic system \eqref{eq:char_x}--\eqref{eq:char_p} is a
\textbf{Hamiltonian system} in the sense of classical mechanics, with $x$ and
$p$ playing the roles of generalized coordinates and momenta, and $H(t,x,p)$
playing the role of the Hamiltonian. Such systems preserve the \textbf{symplectic
form} $\omega = \sum_i dp_i \wedge dx_i$ --- a result known as Liouville's
theorem in its classical form, or more precisely as the preservation of the
symplectic structure by Hamiltonian flows.

This symplectic structure is what guarantees the identity \eqref{eq:bridge}
persists along the flow: the hypersurface $\{p = \nabla_x\phi(t,x)\}$ in
$(x,p)$-space is a \textbf{Lagrangian submanifold} of the symplectic phase
space (a submanifold on which $\omega$ vanishes), and Lagrangian submanifolds
are preserved by Hamiltonian flows. The initial condition $p(T) = \nabla_x G(x)$
defines a Lagrangian submanifold at $t = T$ (the graph of the gradient of a
scalar function is always Lagrangian), and the Hamiltonian flow propagates it
backward in time as a Lagrangian submanifold --- which means $p(t) =
\nabla_x\phi(t, x(t))$ continues to hold for all $t$. Caustics (where $\phi$
loses smoothness) correspond to the Lagrangian submanifold developing folds in
the $x$-projection, which is precisely where the value function becomes
non-differentiable and viscosity solutions are needed.

For the neural network approximation, this geometric picture provides intuition
for why learning $\phi_\theta$ is the right parameterization: a neural network
that approximates a scalar field automatically produces a gradient
$\nabla_x\phi_\theta$ that is the gradient of something --- i.e., it
automatically satisfies the integrability condition $\partial_{x_i} p_j =
\partial_{x_j} p_i$ that characterizes Lagrangian submanifolds. Directly
parameterizing $p_x$ as a vector field would not enforce this condition without
additional architectural constraints, and violating it would break the
PMP--HJB bridge entirely.
\end{physicsinsightbox}


\subsection{From the Gradient Identity to the Implicit Solve}
\label{sec:implicit-solve}

\medskip\hrule\medskip

\subsubsection{Where We Stand}

\begin{tldrbox}
Section~\ref{sec:bridge} established the identity $p_x^\star(t) = \nabla_x \phi(t, z_x^\star(t))$
rigorously. We know that PMP and HJB are the same object, that the costate is
the gradient of the value function along optimal trajectories, and that this
identity is preserved. We
also acknowledged that this identity is classical --- it has been known and used
in the optimal control literature for decades.

What was \textit{not} done before the paper is to exploit this identity inside
a neural network training loop for problems with \textbf{implicit Hamiltonians}.
This step makes precise exactly what changes when the Hamiltonian is implicit,
why the identity alone is insufficient without an implicit solve, and how
the fixed-point iteration fills the gap. This is the last purely mathematical step before
Part IV, where the full neural network architecture and training procedure are
developed.
\end{tldrbox}

\medskip\hrule\medskip

\subsubsection{The Explicit Case Revisited: What the Identity Buys Directly}

In the explicit case (Section~\ref{sec:explicit-case}), the optimal control has a closed-form expression:

\begin{equation}
    u^\star = g(t, z_x, p_x).
\end{equation}

Substituting the identity $p_x = \nabla_x \phi_\theta(t, z_x)$:

\begin{equation}
    u^\star = g(t, z_x, \nabla_x \phi_\theta(t, z_x)).
    \label{eq:explicit_substitution}
\end{equation}

This is immediately computable: given the current state $z_x$ and the network
$\phi_\theta$, one forward pass through $\phi_\theta$ followed by
automatic differentiation gives $\nabla_x \phi_\theta$, and then $g$ maps this
to $u^\star$ in closed form. No iteration, no inner loop. The dynamics can be
integrated forward using a standard ODE solver with $u^\star$ evaluated at each
step via \eqref{eq:explicit_substitution}. This is precisely what prior
neural-network value-function methods do, and it works seamlessly as long as
$g$ is available.

The canonical example from Section~\ref{sec:explicit-case} is $u^\star = -R^{-1} B(t,z_x)^\top
\nabla_x \phi_\theta(t, z_x)$, a simple matrix-vector product applied to the
gradient of the network. The cost of evaluating $u^\star$ at each ODE step is
$O(mn)$ --- one matrix-vector multiply --- on top of the autodiff cost.

\medskip\hrule\medskip

\subsubsection{The Implicit Case: What Breaks}

Now suppose the Hamiltonian is implicit: the first-order condition
$\nabla_u \mathcal{H} = 0$ does not yield a usable closed-form expression for
$u^\star$. We still have the identity $p_x = \nabla_x \phi_\theta$, so the
optimality condition becomes:

\begin{equation}
    \nabla_u \mathcal{H}(t, z_x, \nabla_x \phi_\theta(t, z_x), u^\star) = 0.
    \label{eq:implicit_condition}
\end{equation}

This is a nonlinear equation in $u^\star \in \mathbb{R}^m$ that must be solved
at every point $(t, z_x)$ where the control is needed. The issue is not that
the Hamiltonian $H$ fails to exist; it does. The issue is that evaluating
$\arg\sup_u \mathcal{H}$ repeatedly has no closed form and becomes the
computational bottleneck.

Prior methods that parameterized the value function depended on explicit
evaluation of the control map $g$. Without such a map, they could not evaluate
$u^\star$ cheaply enough inside trajectory rollout and residual evaluation.

The same bottleneck appears in classical grid-based HJB:
evaluating $H(t, x, \nabla_x\phi)$ at each grid point requires solving
$\sup_u \mathcal{H}$, which for an implicit Hamiltonian requires an inner
iterative solve at every grid point at every time step. The neural network
setting does not introduce this difficulty --- it was always there. What the
neural network setting does is make the problem tractable in high dimensions
\textit{if} the implicit solve can be handled efficiently, which is precisely
what the paper achieves.

\medskip\hrule\medskip

\subsubsection{The Implicit Function Theorem as the Key}

\begin{mathrequirementsbox}
The resolution is to treat \eqref{eq:implicit_condition} as a
\textbf{fixed-point equation} and solve it iteratively. Define:

\begin{equation}
    F(u;\, t, z_x, \theta)
    := \nabla_u \mathcal{H}(t, z_x, \nabla_x \phi_\theta(t, z_x), u)
    = 0.
    \label{eq:fixed_point}
\end{equation}

This is a system of $m$ equations in $m$ unknowns $u \in \mathbb{R}^m$,
parameterized by $(t, z_x, \theta)$. The \textbf{implicit function theorem}
(IFT) guarantees that $u^\star$ is a well-defined smooth function of
$(t, z_x, \theta)$ in a neighborhood of any solution, provided the Jacobian
$\partial F / \partial u$ is invertible at the solution:

\begin{equation}
    \frac{\partial F}{\partial u}\bigg|_{u^\star}
    = \nabla^2_{uu} \mathcal{H}(t, z_x, \nabla_x\phi_\theta, u^\star).
    \label{eq:IFT_condition}
\end{equation}

This Jacobian is invertible if and only if $\nabla^2_{uu}\mathcal{H} \prec 0$
--- that is, if $\mathcal{H}$ is \textbf{strictly concave in $u$} at $u^\star$.
This is the same concavity condition identified earlier as necessary
for the implication $u^\star \in \arg\max_u \mathcal{H} \Rightarrow \nabla_u
\mathcal{H} = 0$ to hold and for the maximizer to be unique. Concavity of
$\mathcal{H}$ in $u$ therefore plays three roles simultaneously:

\begin{enumerate}
    \item It guarantees that the first-order condition $\nabla_u \mathcal{H} = 0$
          is both necessary \textit{and} sufficient for a global maximum.
    \item It guarantees the maximizer $u^\star$ is unique.
    \item It guarantees via the IFT that $u^\star(t, z_x, \theta)$ is a
          smooth function of all its arguments, and in particular is
          differentiable with respect to the network parameters $\theta$
          (this step).
\end{enumerate}

The third role is the new one here. Differentiability of $u^\star$ with respect
to $\theta$ is essential for training: the loss function depends on trajectories
generated under $u^\star$, which depends on $\theta$ through
$\nabla_x\phi_\theta$, and backpropagating the training loss requires
$\partial u^\star / \partial \theta$. The IFT provides this derivative
implicitly:

\begin{equation}
    \frac{\partial u^\star}{\partial \theta}
    = -\left[\nabla^2_{uu}\mathcal{H}\right]^{-1}
    \nabla_{u,\theta}^2 \mathcal{H}
    \label{eq:IFT_derivative}
\end{equation}

where $\nabla^2_{u,\theta}\mathcal{H} = \partial^2 \mathcal{H} /
\partial u\, \partial\theta$ captures how the stationarity condition changes
with $\theta$ (through the dependence of $\nabla_x\phi_\theta$ on $\theta$).
This is the exact gradient of $u^\star$ with respect to $\theta$. Computing it
requires solving one linear system involving $\nabla^2_{uu}\mathcal{H}$ per
parameter --- which is expensive and is precisely what
\textbf{Jacobian-Free Backpropagation} avoids. We return to this in Section~\ref{sec:jfb}.
\end{mathrequirementsbox}

\medskip\hrule\medskip

\subsubsection{The Fixed-Point Iteration as the Inner Solve}

\begin{compnumissuesbox}
Given the guarantee from the IFT that $u^\star$ exists and is unique under
strict concavity, we solve \eqref{eq:fixed_point} using \textbf{gradient ascent}
as a fixed-point iteration. Define the operator:

\begin{equation}
    T_\theta(u;\, t, z_x)
    = u + \alpha\, \nabla_u \mathcal{H}(t, z_x, \nabla_x\phi_\theta(t, z_x), u)
    \label{eq:fixed_point_iter}
\end{equation}

for a step size $\alpha > 0$. The fixed point $u^\star$ of $T_\theta$ satisfies
$T_\theta(u^\star) = u^\star$, which is equivalent to $\nabla_u \mathcal{H}
= 0$. Starting from an initial guess $u^{(0)}$ (e.g., zero), iterate:

\begin{equation}
    u^{(k+1)} = T_\theta(u^{(k)}) = u^{(k)}
    + \alpha\, \nabla_u \mathcal{H}(t, z_x, \nabla_x\phi_\theta, u^{(k)})
    \label{eq:gd_iter}
\end{equation}

until $\|u^{(k+1)} - u^{(k)}\| < \varepsilon_{\text{tol}}$ (or a maximum of $K$
steps). Convergence is guaranteed by the paper's \textbf{Assumption~1}:
$T_\theta$ is $\gamma$-contractive with $\gamma < 1$ for sufficiently small
$\alpha$, so the iterates converge \textbf{linearly} to $u^\star$ at rate
$\gamma$:
$\|u^{(k)} - u^\star\| \leq \gamma^k \|u^{(0)} - u^\star\|$.
The paper uses $K = 200$ with tolerance $10^{-3}$ or $10^{-4}$ depending on
the problem (see the experimental setup sections).

Note that this iteration requires only the \textbf{first derivative}
$\nabla_u \mathcal{H}$, not the Hessian $\nabla^2_{uu}\mathcal{H}$. The code
initializes $u^{(0)} = 0$ at every time step (cold start); no warm-starting
is used.

\textbf{Classical alternative: Newton's method.} A classical approach to solving
$\nabla_u \mathcal{H} = 0$ is Newton's method:
$u^{(k+1)} = u^{(k)} - [\nabla^2_{uu}\mathcal{H}]^{-1}\, \nabla_u \mathcal{H}$.
This converges quadratically but requires computing and inverting the $m \times m$
Hessian $\nabla^2_{uu}\mathcal{H}$ at each step. The paper does \textit{not} use
Newton's method for the forward solve --- it uses the gradient ascent operator
\eqref{eq:gd_iter} instead.

\textbf{Where the Hessian does appear.} The Hessian $\nabla^2_{uu}\mathcal{H}$
appears legitimately in the \textit{exact gradient} of $u^\star$ with respect to
$\theta$ via the IFT: $\partial u^\star / \partial\theta = -[\nabla^2_{uu}
\mathcal{H}]^{-1}\, \nabla^2_{u,\theta}\mathcal{H}$ (equation
\eqref{eq:IFT_derivative}). This is part of Term~II --- the expensive gradient
term that JFB avoids computing. The Hessian is needed for the \textit{backward
pass} (exact differentiation), not the \textit{forward solve}.

\textbf{Dependence on $\nabla_x\phi_\theta$.} The iterate
\eqref{eq:gd_iter} depends on $\nabla_x\phi_\theta(t, z_x)$, which is obtained
by differentiating the network with respect to its input $x$. This is a standard
autodiff operation. Crucially, the fixed-point iteration itself is \textit{not}
differentiated through during training --- this is the stop-gradient operation
that defines JFB and is justified in Section~\ref{sec:jfb}.
\end{compnumissuesbox}

\medskip\hrule\medskip

\subsubsection{The Complete Control Evaluation Pipeline}

\begin{tldrbox}
We can now state the complete procedure for evaluating $u^\star(t, z_x;\theta)$
at a single point during trajectory generation, combining the PMP--HJB bridge
with the implicit solve:

\medskip

\begin{enumerate}
    \setlength{\itemsep}{0.6em}

    \item \textbf{Forward pass:} evaluate $\phi_\theta(t, z_x)$ via the neural
          network. Cost: one forward pass, $O(\text{parameters})$.

    \item \textbf{Autodiff:} compute $\nabla_x\phi_\theta(t, z_x)$ by
          differentiating the forward pass with respect to the input $x$.
          Cost: one backward pass through the network, same order as forward.

    \item \textbf{Set costate:} $p_x \leftarrow \nabla_x\phi_\theta(t, z_x)$
          via the bridge identity of Section~\ref{sec:bridge}.

    \item \textbf{Fixed-point solve:} iterate $u \leftarrow u + \alpha\,
          \nabla_u\mathcal{H}(t, z_x, p_x, u)$ until
          $\|u_{\text{new}} - u_{\text{old}}\| < \varepsilon_{\text{tol}}$.
          Convergence guaranteed by $\gamma$-contractivity of $T_\theta$
          (Assumption~1). Cost: $K$ iterations, each $O(m)$ for the
          gradient evaluation.

    \item \textbf{Return $u^\star$} to the ODE integrator for the dynamics step
          $\dot{z}_x = f(t, z_x, u^\star)$.
\end{enumerate}

\medskip

This pipeline replaces the single closed-form evaluation $u^\star =
g(t, z_x, \nabla_x\phi_\theta)$ of the explicit case with a short iterative
loop, at a cost that is $O(K \cdot m)$ per control evaluation (one gradient
evaluation per step). For the physical problems the paper considers --- space
shuttle ($m = 1$), bicycle ($m = 2$) --- this cost is negligible.
\end{tldrbox}

\medskip\hrule\medskip

\subsubsection{What Remains: The Training Problem}

\begin{tldrbox}
The pipeline above describes how to \textit{evaluate} $u^\star$ given a fixed
network $\phi_\theta$. The remaining question is how to \textit{train}
$\phi_\theta$ --- how to choose $\theta$ so that $\phi_\theta \approx \phi$,
the true value function.

Training requires:
\begin{enumerate}
    \item A \textbf{loss function} --- the control objective $J(\theta)$,
          the expected cost of trajectories rolled out under the current policy.
    \item A source of \textbf{training data} --- trajectories generated by
          simulating the controlled dynamics under the current
          $u^\star(\cdot;\theta)$.
    \item A \textbf{gradient computation} --- the derivative of $J$ with
          respect to $\theta$, used to update $\theta$ via gradient descent.
\end{enumerate}

The first two are straightforward. The third is where the \textbf{temporal
coupling} problem arises: the training trajectories depend on $\theta$ (because
$u^\star$ depends on $\phi_\theta$), so the gradient of the loss with respect
to $\theta$ has two terms --- one explicit (through the cost at fixed trajectory
points) and one implicit (through how the trajectories change with $\theta$).
The implicit term requires differentiating through the ODE solver and through
the fixed-point iterations, which is the computational bottleneck that
\textbf{Jacobian-Free Backpropagation} resolves.

The remaining sections develop the training procedure in full, culminating in the
JFB theorem and its proof. Section~\ref{sec:nn-param} introduces the neural network
architecture and the training objective. Section~\ref{sec:training-temporal-coupling}
defines the training loop, derives the gradient split, and makes the temporal
coupling explicit. Section~\ref{sec:jfb-descent} proves the descent direction
theorem. Section~\ref{sec:jfb} presents the complete JFB algorithm.
\end{tldrbox}

\pagebreak

\section{Part IV --- Neural Network Approach and Jacobian-Free Backpropagation}

\paragraph{Overview of Part IV.}
Before developing the formalism, we state the complete narrative of this part
in plain terms so that each subsequent section can be read against it.

\textbf{What we are optimizing.} The training loss is the control objective
$J(\theta) = \mathbb{E}_{x\sim\rho}[\int_0^T L\,dt + \alpha_G\, G(z(T))]$
--- the expected cost of trajectories rolled out under the current policy.
The HJB residual is \emph{never} in the training loss; it is computed as a
diagnostic metric only, with zero weight in the optimizer.

\textbf{Why this creates a hard gradient problem.} Computing $dJ/d\theta$
requires differentiating through the entire trajectory rollout. At each of
$N_t$ time steps, the control $u^\star_k$ was computed from
$\nabla_z\phi_\theta(t_k, z_k)$, and the state $z_k$ itself was steered
there by all preceding controls. This temporal coupling makes the exact
gradient intractable for any network with more than a few thousand parameters.

\textbf{What JFB does.} The gradient splits into two terms: Term~I, which
differentiates through $\phi_\theta$ at fixed trajectory points (cheap,
$O(P)$ per step), and Term~II, which differentiates through the trajectory
itself (expensive, requires the full sensitivity ODE). JFB keeps Term~I and
discards Term~II by treating trajectory states as constants via
\texttt{stop\_grad}.

\textbf{Why this is justified.} Theorem~1 proves that Term~I alone is a
descent direction for $J(\theta)$ --- its inner product with the true
gradient is strictly positive --- under four assumptions about contractivity,
conditioning, and regularity. The proof is algebraic and finite-time; it
works by establishing a pointwise inner product bound
$\langle v_\theta(t), w_\theta(t)\rangle > 0$ and lifting it to the
time-integrated gradient via Lemma~4.

\textbf{Training versus inference.} At inference, the fixed-point iteration
runs to convergence with no gradient tracking. During training, the iteration
first converges in a \texttt{no\_grad} block, then \texttt{tracked\_iters}
additional applications of $T_\theta$ are executed with gradient tracking,
injecting the gradient signal into $\theta$ through $\phi_\theta$.

\subsection{Neural Network Parameterization of the Value Function}
\label{sec:nn-param}

\medskip\hrule\medskip

\subsubsection{Why Neural Networks}

\begin{extrainfobox}
Classical grid-based representations of $\phi(t,x)$ require $O(N^n)$ storage
--- intractable for $n \gtrsim 6$ (Section~\ref{sec:curse}). A neural network
$\phi_\theta : [0,T] \times \mathbb{R}^n \to \mathbb{R}$ with parameters
$\theta \in \mathbb{R}^P$ sidesteps this: evaluation costs $O(P)$ regardless
of $n$, and the gradient $\nabla_x\phi_\theta$ needed for the costate
substitution and the fixed-point iteration is obtained for free via autodiff
through the forward pass.
\end{extrainfobox}

\medskip\hrule\medskip

\subsubsection{The Network Architecture}

The paper parameterizes $\phi$ as a deep network
$\phi_\theta : [0,T] \times \mathbb{R}^n \to \mathbb{R}$ with parameters
$\theta \in \mathbb{R}^P$. Since we need $\nabla_x\phi_\theta$ for the
costate substitution and fixed-point iteration, the network must be at least
$C^1$ in its inputs. This rules out ReLU in favor of smooth activations;
the paper uses $\tanh$ throughout.

\textbf{Enforcing the terminal condition.} The terminal condition $\phi(T,x) =
G(x)$ must be satisfied by $\phi_\theta$. There are two approaches. The first
is \textbf{soft enforcement}: include a terminal penalty term in the loss
function that penalizes $\|\phi_\theta(T,x) - G(x)\|^2$, allowing the network
to learn to satisfy the condition approximately. The second is \textbf{hard
enforcement}: reparameterize the network so that the terminal condition is
satisfied by construction. A standard hard enforcement is:

\begin{equation}
    \phi_\theta(t,x) = G(x) + (T - t)\, \mathcal{N}_\theta(t,x)
    \label{eq:hard_terminal}
\end{equation}

which satisfies $\phi_\theta(T,x) = G(x)$ exactly for any $\theta$, since
the factor $(T-t)$ vanishes at $t=T$. The paper uses hard enforcement via
\eqref{eq:hard_terminal}, which removes the terminal condition from the
optimization entirely and reduces the effective degrees of freedom the network
must learn.

\medskip\hrule\medskip

\subsubsection{The Training Objective}

\begin{quote}
\textbf{Important.} The HJB residual described in this section is \emph{not}
the training loss. It is a diagnostic quantity computed during training to
verify that $\phi_\theta$ is converging toward the true value function, but
its weight in the optimizer is set to zero ($\alpha_{\text{HJB}} = [0,0]$,
$\alpha_{\text{adj}} = [0,0]$). The actual training loss is the control
objective $J(\theta)$, defined below and used in the training loop of
Section~\ref{sec:training-temporal-coupling}. The fact that the HJB residual
decreases during training even though it is never minimized is one of the
paper's key empirical findings: minimizing $J(\theta)$ end-to-end implicitly
pushes $\phi_\theta$ toward satisfying the HJB equation, by virtue of the
PMP--HJB bridge established in Part~III.

\end{quote}

The paper trains $\phi_\theta$ by directly minimizing the \textbf{control
objective} --- the expected cost of trajectories rolled out under the current
policy. For a single initial condition $x$, define:

\begin{equation}
    J_x(\theta) = \int_0^T L(t, z_x, u^\star_\theta)\, dt
    + \alpha_G \cdot G(z_x(T))
    \label{eq:control_objective}
\end{equation}

where $z_x$ is the state trajectory obtained by integrating the controlled
dynamics $\dot{z}_x = f(t, z_x, u^\star_\theta)$ forward from initial
condition $z_x(0) = x$, and $u^\star_\theta$ is the optimal control recovered
via the fixed-point iteration of Section~\ref{sec:implicit-solve} at each time
step. The training objective is the expectation over a distribution $\rho$ of
initial conditions:

\begin{equation}
    \min_\theta\; \mathbb{E}_{x \sim \rho}\bigl[J_x(\theta)\bigr]
    \label{eq:training_objective}
\end{equation}

This choice of training loss has a direct consequence for the gradient
computation. Because $J_x(\theta)$ depends on $\theta$ through entire
rolled-out trajectories --- and each point on those trajectories was generated
using $u^\star_\theta$, which itself depends on $\phi_\theta$ --- the gradient
$dJ/d\theta$ accumulates coupling across all time steps. This is the
\textbf{temporal coupling} problem developed in
Section~\ref{sec:training-temporal-coupling}.

\medskip\hrule\medskip

\subsubsection{Why Parameterizing $\phi_\theta$ Rather Than $u_\theta$}

\begin{extrainfobox}
An alternative approach --- used in reinforcement learning and direct policy
optimization --- is to parameterize the control policy directly:
$u_\theta(t,x) \approx u^\star(t,x)$. This avoids the need to solve any PDE
and produces a feedback law immediately. Why does the paper instead parameterize
the value function?

There are three reasons, each progressively more important.

\textbf{First}, directly parameterizing $u_\theta$ does not leverage the
physical structure of the problem. The value function $\phi$ satisfies the HJB
equation --- a hard constraint with physical content. Training $\phi_\theta$ to
satisfy this constraint injects physical knowledge into the learning process and
acts as a strong regularizer. A directly parameterized policy $u_\theta$ has no
such structural constraint and must learn the correct behavior entirely from
data.

\textbf{Second}, the gradient $\nabla_x\phi_\theta$ provides the costate via
the bridge identity, which in turn provides the correct direction for the
implicit solve. A directly parameterized $u_\theta$ cannot use this bridge ---
it has no costate, no Hamiltonian structure, and no way to handle implicit
Hamiltonians via the fixed-point iteration. The implicit solve \textit{requires} the
costate, which \textit{requires} the value function parameterization.

\textbf{Third}, the value function encodes the solution for all initial
conditions simultaneously. A policy $u_\theta(t,x)$ does too in principle, but
without the HJB constraint there is no guarantee that $u_\theta$ is consistent
across different initial conditions in the way that $\nabla_x\phi_\theta$
automatically is (by the Lagrangian submanifold structure of Section~\ref{subsec:lagrangian-submanifold}).

These three reasons together explain why the paper's architecture is
$\phi_\theta$ rather than $u_\theta$, and why this choice is not merely
conventional but is structurally necessary for the implicit Hamiltonian approach.
\end{extrainfobox}

\medskip\hrule\medskip

\subsubsection{Prior Work and the Explicit Hamiltonian Requirement}

\begin{extrainfobox}
Several prior works parameterize the value function as a neural network and
train it to satisfy the HJB equation. The most directly related are methods
based on the \textbf{Hopf formula} and \textbf{data-driven HJB solvers} such as
the work of Darbon and Osher, and the neural network approaches of Nakamura-Zimmerer
et al.\ and Onken et al. All of these methods share a common requirement: they
need a closed-form expression for $u^\star = g(t, z_x, \nabla_x\phi_\theta)$ to:

\begin{enumerate}
    \item Generate training trajectories by integrating the dynamics forward
          under $u^\star = g(t, z_x, \nabla_x\phi_\theta)$ at each time step.
    \item Evaluate the running cost $L(t, z_x, u^\star)$ and terminal cost
          $G(z_x(T))$ along those trajectories.
    \item Backpropagate $dJ/d\theta$ through the trajectory rollout to update
          $\theta$ via gradient descent.
\end{enumerate}

Without $g$, step~1 breaks: you cannot integrate the trajectory forward
because you cannot evaluate the control at each step. The paper resolves this
by replacing $g$ with the fixed-point iteration of
Section~\ref{sec:implicit-solve}, and replacing exact backpropagation through
the rollout with JFB. This is the precise gap the paper fills in the landscape
of neural HJB methods.
\end{extrainfobox}

\medskip

\hrule

\subsection{End-to-End Training and Temporal Coupling}
\label{sec:training-temporal-coupling}

\medskip\hrule\medskip

\subsubsection{The Training Loop}

We now describe the complete training procedure and identify precisely where
the temporal coupling arises and why it makes exact gradient computation
expensive.

At each training iteration, a batch of initial conditions
$\{x_i\}_{i=1}^{B}$ is sampled from the distribution $\rho$. For each
$x_i$, a trajectory is generated by integrating the controlled dynamics
forward:

\begin{equation}
    \dot{z}_{x_i}(t) = f(t, z_{x_i}(t), u^\star(t, z_{x_i}(t);\theta)),
    \qquad z_{x_i}(0) = x_i
    \label{eq:controlled_dynamics}
\end{equation}

using forward Euler with step size $h = T/N_t$. At each of the $N_t$ time
steps, the control $u^\star_\theta$ is evaluated via the pipeline of
Section~\ref{sec:implicit-solve}: autodiff gives $\nabla_z\phi_\theta$, then
fixed-point iteration converges to $u^\star$. The running cost is accumulated
as $\sum_k h \cdot L(t_k, z_k, u^\star_k)$, and the terminal cost
$G(z(T))$ is evaluated at the end of the rollout.

The total training loss over the batch is:

\begin{equation}
    \mathcal{L}(\theta)
    = \frac{1}{B} \sum_{i=1}^{B}
    \left[\alpha_L \sum_{k=0}^{N_t-1} h \cdot L(t_k, z_{x_i}(t_k), u^\star_k)
    + \alpha_G \cdot G(z_{x_i}(T))\right].
    \label{eq:training_loss}
\end{equation}

Parameters are updated by gradient descent: $\theta \leftarrow \theta - \eta
\nabla_\theta \mathcal{L}(\theta)$ for learning rate $\eta > 0$. The entire
procedure is differentiable in principle --- every operation from $\theta$
through the network, through the fixed-point solve, through the ODE
integration, to the cost is differentiable. The question is whether computing
$\nabla_\theta\mathcal{L}$ exactly is tractable.

Note that the HJB residual $\mathcal{R}(t,z;\theta)$ and adjoint residual are
also computed along each trajectory and logged, but with zero weight in
$\mathcal{L}(\theta)$. They serve as diagnostic metrics only.

\medskip\hrule\medskip

\subsubsection{The Cost of Term II: Why Temporal Coupling Is Expensive}

To understand the cost, consider the sensitivity of the trajectory with respect
to $\theta$. Differentiating the dynamics \eqref{eq:controlled_dynamics} gives
the \textbf{sensitivity ODE}:

\begin{equation}
    \frac{d}{dt}\frac{dz_{x_i}}{d\theta}
    = \frac{\partial f}{\partial z}
    \frac{dz_{x_i}}{d\theta}
    + \frac{\partial f}{\partial u^\star}
    \frac{du^\star}{d\theta}
    \label{eq:sensitivity_ode}
\end{equation}

with initial condition $dz_{x_i}(0)/d\theta = 0$. The term
$du^\star/d\theta$ requires the IFT formula \eqref{eq:IFT_derivative} from
Section~\ref{sec:implicit-solve}:

\begin{equation}
    \frac{du^\star}{d\theta}
    = -\left[\nabla^2_{uu}\mathcal{H}\right]^{-1}
    \frac{\partial}{\partial\theta}
    \nabla_u\mathcal{H}(t, z, \nabla_x\phi_\theta, u^\star)
    \label{eq:dustar_dtheta}
\end{equation}

which requires differentiating $\nabla_x\phi_\theta$ with respect to $\theta$
--- a second-order autodiff operation --- and solving an $m \times m$ linear
system at every time step.

This is a linear
ODE of dimension $n \times P$ (state dimension times number of parameters)
that must be integrated alongside the main trajectory ODE. Its solution at time
$t_j$ is:

\begin{equation}
    \frac{dz_{x_i}(t_j)}{d\theta}
    = \int_0^{t_j} \Phi(t_j, s)
    \frac{\partial f}{\partial u^\star}(s)
    \frac{du^\star}{d\theta}(s)\, ds
    \label{eq:sensitivity_integral}
\end{equation}

where $\Phi(t_j, s)$ is the \textbf{state transition matrix} of the linearized
dynamics, satisfying $\partial_t\Phi(t,s) = (\partial f/\partial z)\Phi(t,s)$
with $\Phi(s,s) = I_n$. Each column of $\Phi$ requires integrating an $n$-
dimensional linear ODE, and there are $P$ columns --- one per network parameter.
Since $P$ can be $10^4$ to $10^6$ for typical neural networks, integrating the
full sensitivity ODE is completely intractable.

The equivalent \textbf{adjoint method} (the continuous-time analogue of
backpropagation through time) computes $\nabla_\theta\mathcal{L}$ by integrating
a single adjoint ODE backward in time of dimension $n$, followed by a single
forward pass. This reduces the memory cost from $O(nP)$ to $O(n)$, but still
requires:

\begin{enumerate}
    \item Storing the full forward trajectory $z_{x_i}(t)$ for $t \in [0,T]$
          (or recomputing it via checkpointing) --- memory cost $O(M \cdot n)$
          per trajectory.
    \item Integrating the adjoint ODE backward, which requires evaluating
          $\partial f/\partial u^\star$ and $du^\star/d\theta$ at every time
          step --- each requiring the IFT solve \eqref{eq:dustar_dtheta}.
    \item The IFT solve itself requires $\nabla^2_{uu}\mathcal{H}$ and a
          second-order autodiff through $\phi_\theta$, which is typically
          $O(P^2)$ in memory and $O(P^2 n)$ in computation for dense networks.
\end{enumerate}

Even with the adjoint method, the per-iteration cost of computing the exact
gradient is dominated by the second-order autodiff required by
\eqref{eq:dustar_dtheta}. For networks with $P = 10^5$ parameters, this is
simply not feasible.

For implicit Hamiltonians, the coupling has an additional layer: $u^\star$
depends on $\theta$ not only through $\phi_\theta$ at the trajectory points,
but also through the fixed-point solve via \eqref{eq:dustar_dtheta}. The JFB
approach of Section~\ref{sec:jfb} addresses both layers simultaneously.

\medskip\hrule\medskip

\subsubsection{The HJB Residual as a Diagnostic}

\begin{extrainfobox}
Although the training loss is the control objective $J(\theta)$, the HJB
residual $\mathcal{R}(t,z;\theta) = -\partial_t\phi_\theta +
H(t,z,\nabla_x\phi_\theta)$ is computed alongside each trajectory during
training and logged as a diagnostic. Its weight in the optimizer is zero.

This separation is deliberate. The paper's key empirical observation is that
$\|\mathcal{R}\|$ decreases during training even though it is never directly
minimized. This confirms that minimizing $J(\theta)$ end-to-end implicitly
drives $\phi_\theta$ toward the true value function --- a consequence of the
PMP--HJB bridge: if $\phi_\theta$ produces trajectories with low cost, it must
be a good approximation to the true value function, which satisfies the HJB
equation exactly.

This also explains why the temporal coupling problem is unavoidable: the
training signal comes from trajectory cost, so the training states are
trajectory states, which depend on $\theta$. A PINN-style approach that
directly minimized the HJB residual at fixed sample points would have no
temporal coupling, but would lose the end-to-end structure and would require
evaluating $H(t,x,\nabla_x\phi_\theta)$ --- the implicit supremum --- at
every sample point as part of the loss itself, not just during trajectory
generation.
\end{extrainfobox}

\medskip\hrule\medskip

\subsubsection{The Two-Term Gradient Split}
\label{sec:full-gradient}

For a single trajectory starting from $x$, discretize time as $0 = t_0 <
t_1 < \cdots < t_M = T$ with step size $h = T/M$. Write
$z_k = z_x(t_k;\theta)$ and $u_k^\star = u^\star(t_k, z_k;\theta)$. The
discrete control objective is:

\begin{equation}
    \mathcal{L}(\theta)
    = \sum_{k=0}^{M-1} h \cdot L(t_k, z_k, u^\star_k(\theta))
    + \alpha_G \cdot G(z_M(\theta)).
    \label{eq:discrete_loss}
\end{equation}

The full gradient is:

\begin{equation}
    \frac{d\mathcal{L}}{d\theta}
    = \underbrace{
    \sum_{k=0}^{M-1} h \cdot
    \frac{\partial L(t_k, z_k, u^\star_k)}{\partial\theta}
    \bigg|_{z_k\,\text{fixed}}
    + \alpha_G \cdot \frac{\partial G(z_M)}{\partial\theta}
    \bigg|_{z_M\,\text{fixed}}
    }_{\text{Term I: } \widetilde{\nabla}_\theta\mathcal{L}}
    +\; \underbrace{
    \sum_{k=0}^{M} \nabla_{z_k} C_k \cdot
    \frac{dz_k}{d\theta}}_{\text{Term II: trajectory coupling}}
    \label{eq:full_gradient}
\end{equation}

where $\nabla_{z_k} C_k$ collects the cost gradients with respect to the
state at $t_k$ (running cost at interior steps, terminal cost at $t_M$).
Term~I is the \textbf{JFB gradient}
$\widetilde{\nabla}_\theta\mathcal{L}$ --- the approximation the paper uses.
Term~II is the trajectory coupling term that makes exact computation
intractable.

\medskip\hrule\medskip

\subsubsection{The Trajectory Jacobian}

The quantity $dz_k/d\theta \in \mathbb{R}^{n \times P}$ in Term II is the
\textbf{trajectory Jacobian} --- how the state at time $t_k$ changes with the
network parameters. By the discrete chain rule through the ODE integrator:

\begin{equation}
    \frac{dz_k}{d\theta}
    = \sum_{j=0}^{k-1}
    \left(\prod_{\ell=j+1}^{k-1} \frac{\partial z_{\ell+1}}{\partial z_\ell}
    \right)
    \frac{\partial z_{j+1}}{\partial u_j^\star}
    \frac{du_j^\star}{d\theta}
    \label{eq:traj_jacobian}
\end{equation}

where $\partial z_{\ell+1}/\partial z_\ell$ is the Jacobian of one ODE step
with respect to the state (an $n \times n$ matrix), and
$du_j^\star/d\theta$ is the IFT derivative \eqref{eq:dustar_dtheta}. The
product of ODE step Jacobians $\prod_\ell \partial z_{\ell+1}/\partial z_\ell$
is the discrete-time state transition matrix $\Phi(t_k, t_j)$, whose operator
norm grows as $e^{\lambda_{\max}(t_k - t_j)}$ for unstable systems --- the
same exponential growth that made the PMP shooting problem ill-conditioned.

For implicit Hamiltonians, each $du_j^\star/d\theta$ requires the IFT formula:

\begin{equation}
    \frac{du_j^\star}{d\theta}
    = -\left[\nabla^2_{uu}\mathcal{H}_j\right]^{-1}
    \nabla_{u,\theta}^2\mathcal{H}_j
    \in \mathbb{R}^{m \times P}
\end{equation}

where $\nabla_{u,\theta}^2\mathcal{H}_j = \partial_u\partial_\theta
[\langle\nabla_x\phi_\theta(t_j,z_j), f\rangle + L]$ is a mixed
second-order derivative requiring one Jacobian-vector product through the
network per component of $u$ --- $O(mP)$ operations total per time step,
and $O(mPM)$ over the full trajectory. For $P = 10^5$, $m = 2$, $M = 100$,
this is $O(2 \times 10^7)$ operations per trajectory, per iteration, before
even accounting for the matrix inversion. Combined across a batch of $B$
trajectories, exact gradient computation is three to four orders of magnitude
more expensive than Term I alone.

\medskip\hrule\medskip

\subsubsection{Why Term II Cannot Simply Be Ignored Without Justification}

Dropping Term II and using only Term I as the gradient update is an
approximation. Without justification, this could lead to:

\begin{enumerate}
    \item \textbf{Wrong descent direction:} if Term II is large and points
          opposite to Term I, the JFB update $-\widetilde{\nabla}_\theta
          \mathcal{L}$ could increase the loss rather than decrease it.
    \item \textbf{Divergence:} repeated ascent steps would cause training
          to diverge entirely.
    \item \textbf{Bias:} even if convergence occurs, the fixed point of
          the JFB update might not correspond to a true critical point of
          $\mathcal{L}$.
\end{enumerate}

None of these pathologies occur under the conditions the paper establishes,
because Term I and the full gradient make an angle of less than $90^{\circ}$ ---
guaranteeing that Term I is a descent direction. This is the content of
the descent-direction theorem in Section~\ref{sec:jfb-descent}.

\medskip

\hrule

\subsection{Jacobian-Free Backpropagation --- The Descent Direction Theorem}
\label{sec:jfb-descent}

\medskip\hrule\medskip

\subsubsection{Statement and Proof}

\begin{mathrequirementsbox}
\label{thm:jfb-descent}
The paper's central theoretical result is that the JFB gradient
$\widetilde{\nabla}_\theta\mathcal{L}$ --- Term~I of
\eqref{eq:full_gradient} --- is a descent direction for
$\mathcal{L}(\theta)$, meaning the update
$\theta \leftarrow \theta - \eta\,\widetilde{\nabla}_\theta\mathcal{L}$
decreases $\mathcal{L}$ for sufficiently small $\eta$. Writing
$\nabla_\theta\mathcal{L} = \widetilde{\nabla}_\theta\mathcal{L} + E(\theta)$
where $E(\theta)$ is Term~II, a sufficient condition by Cauchy--Schwarz is:

\begin{equation}
    \|E(\theta)\| < \|\widetilde{\nabla}_\theta\mathcal{L}\|
    \label{eq:norm_condition}
\end{equation}

The paper proves a stronger result: the inner product is strictly positive
without bounding $\|E\|$ directly.

From \eqref{eq:full_gradient} and \eqref{eq:traj_jacobian}, Term~II satisfies:

\begin{equation}
    \|E(\theta)\|
    \leq \sum_{k=0}^{M}
    \|\nabla_{z_k} C_k\|\,
    \left\|\frac{dz_k}{d\theta}\right\|.
\end{equation}

The trajectory Jacobian $\|dz_k/d\theta\|$ grows as $e^{\lambda_{\max}t_k}$
in the worst case. However, the cost gradients $\|\nabla_{z_k} C_k\|$ remain
bounded, and for a well-trained network the trajectory is nearly optimal so
Term~II is small relative to Term~I.

The paper's actual proof does not bound $\|E(\theta)\|$ directly via a
discretization argument. Instead, it works at the level of the
\textbf{time-integrated gradient vectors}. Define:
\begin{align}
    v_\theta(t) &= \left(\frac{du^\star_\theta}{d\theta}\right)^{\!\top}
    \bigl(\nabla_u L + \nabla_u f^\top p_x\bigr)
    &&\text{(true gradient integrand)}, \\
    w_\theta(t) &= \left(\frac{\partial T_\theta}{\partial\theta}\right)^{\!\top}
    \bigl(\nabla_u L + \nabla_u f^\top p_x\bigr)
    &&\text{(JFB gradient integrand)},
\end{align}
so that $d\mathcal{J}/d\theta = \int v_\theta\, dt$ and
$d\widetilde{\mathcal{J}}/d\theta = \int w_\theta\, dt$.

The descent condition
$\langle d\widetilde{\mathcal{J}}/d\theta,\; d\mathcal{J}/d\theta\rangle > 0$
is equivalent to
$\bigl\langle \int w_\theta\, dt,\; \int v_\theta\, dt \bigr\rangle > 0$.

\textbf{Lemma~2} (under Assumptions~1--3) establishes the \textit{pointwise} bound:
\begin{equation}
    \langle v_\theta(t),\, w_\theta(t) \rangle
    \;\geq\; \|M_\theta\, v_\theta\|^2\,(\lambda_- - \gamma\,\lambda_+)
    \;=:\; \delta^2 > 0
    \label{eq:pointwise_bound}
\end{equation}
for all $t, z, u, \theta$, where $\lambda_-$ and $\lambda_+$ are the
uniform eigenvalue bounds on $(M_\theta M_\theta^\top)^{-1}$ from Lemma~1,
and $\gamma$ is the contractivity constant from Assumption~1. The condition
$\lambda_- - \gamma\,\lambda_+ > 0$ is equivalent to the condition number
$\kappa(M_\theta M_\theta^\top) < 1/\gamma$, which is enforced by
Assumption~2.

\textbf{Lemma~4} then lifts this pointwise bound to the integrated vectors:
if $\langle v_\theta(t),\, w_\theta(t)\rangle \geq \delta^2 > 0$ for all
$t$, and $v_\theta, w_\theta$ satisfy the $L^2$-boundedness and fluctuation
conditions of Assumption~4, then
$\bigl\langle \int v_\theta\, dt,\; \int w_\theta\, dt \bigr\rangle > 0$.
The proof uses the identity:
\begin{equation}
    \langle I_1, I_2 \rangle
    = T \int \langle v_\theta, w_\theta\rangle\, dt
    - T \int \langle v_\theta - C_v,\; w_\theta - C_w\rangle\, dt
\end{equation}
and Assumption~4 ensures the fluctuation term is dominated by the
pointwise inner product term.

\textbf{Theorem~1} follows immediately: $-d\widetilde{\mathcal{J}}/d\theta$
is a descent direction for $\mathcal{J}_x(\theta)$ because its inner product
with the true gradient $-d\mathcal{J}/d\theta$ is positive. The proof is
finite-time and algebraic --- there is no discretization $\Delta t$ and no
asymptotic argument.

\textbf{Connection to prior JFB results.} The name ``Jacobian-Free
Backpropagation'' originates in implicit neural networks (Fung et al., 2022),
where it refers to the same stop-gradient idea applied to static
fixed-point layers. The descent direction proof in the optimal control
context is new: the error term $E(\theta)$ involves the trajectory
Jacobian --- a dynamical systems quantity with no analogue in the static
setting --- so the pointwise inner product bound \eqref{eq:pointwise_bound}
and its lift to the integrated gradient (Theorem~1) are a genuine
contribution of the paper.
\end{mathrequirementsbox}

\medskip\hrule\medskip

\subsubsection{Implications for Training}

\begin{mathrequirementsbox}
The descent direction result has two concrete implications for the training
procedure.

\textbf{Convergence.} Gradient descent with step size $\eta$ satisfying the
Armijo condition decreases $\mathcal{L}$ at every iteration as long as
\eqref{eq:norm_condition} holds. Since $\mathcal{L} \geq 0$, the sequence
$\{\mathcal{L}(\theta^{(k)})\}$ converges. The limit point satisfies
$\widetilde{\nabla}_\theta\mathcal{L} = 0$, which under
\eqref{eq:norm_condition} implies $\nabla_\theta\mathcal{L} = 0$ --- a true
critical point of $\mathcal{L}$.

\textbf{Fixed points.} The fixed points of JFB gradient descent coincide
with the critical points of $\mathcal{L}$ under condition
\eqref{eq:norm_condition}. The JFB approximation therefore does not
introduce spurious fixed points --- any parameter $\theta$ where JFB
training stalls is also a stall point of exact gradient descent.
\end{mathrequirementsbox}


\subsection{Jacobian-Free Backpropagation --- The Complete Algorithm}
\label{sec:jfb}

\medskip\hrule\medskip

\subsubsection{Defining JFB Precisely}

The cost at each time step depends on $\theta$ through two channels:
(1)~\textbf{explicitly}, through $u^\star_k(\theta)$ which depends on
$\nabla_z\phi_\theta$ via $T_\theta$; and
(2)~\textbf{implicitly}, through the trajectory state $z_k$ itself, which
depends on $\theta$ through all preceding controls. JFB blocks channel~2 by
applying \texttt{stop\_grad} to the trajectory states, keeping only channel~1.
The JFB loss is:

\begin{equation}
    \widetilde{\mathcal{L}}(\theta)
    = \frac{1}{B}\sum_{i=1}^{B}
    \left[\alpha_L \sum_{k=0}^{M-1}
    h \cdot L(t_k, \bar{z}_{x_i}(t_k), u^\star_k(\theta))
    + \alpha_G \cdot G(\bar{z}_{x_i}(T))\right]
    \label{eq:jfb_loss}
\end{equation}

where $\bar{z}_{x_i}(t_k) = \texttt{stop\_grad}(z_{x_i}(t_k;\theta))$.
The stop-gradient applies at two levels: trajectory states $\bar{z}_k$ are
detached (blocking channel~2), and the convergence phase of the fixed-point
iteration is not differentiated through --- only the final $K'$ tracked
applications of $T_\theta$ carry gradient. This is Term~I of
\eqref{eq:full_gradient} exactly.

\subsubsection{Complete algorithm}


\begin{algorithm}[t]
\caption{End-to-end training with implicit Hamiltonians and Jacobian-Free Backpropagation (JFB)}
\label{alg:jfb-training}
\begin{algorithmic}[1]
\Require Dynamics $f$, running cost $L$, terminal cost $G$, initial-condition distribution $\rho$,
         learning rate $\eta$, batch size $B$, number of time steps $N_t$, step size $h=T/N_t$,
         fixed-point tolerance $\varepsilon_{\mathrm{tol}}$, ascent step size $\alpha$,
         max fixed-point iterations $K$, tracked iterations $K'$, cost weights $\alpha_L$, $\alpha_G$.
\Ensure Trained value function network $\phi_{\theta^\star}$.

\State Initialize parameters $\theta$.
\Repeat
  \State Sample initial conditions $\{x_i\}_{i=1}^{B} \sim \rho$.
  \For{$i = 1$ to $B$}
    \State $z \gets x_i$; \quad $\mathrm{running\_cost}_i \gets 0$.
    \For{$k = 0$ to $N_t-1$}
      \State $t \gets k h$.
      \State $p_x \gets \nabla_x\phi_\theta(t,z)$ \Comment{autodiff; detach $z$ in code}
      \State $u \gets \mathbf{0}$
      \For{$\ell = 1$ to $K$} \Comment{converge to $u^\star$; no gradient tracking}
        \State $u_{\mathrm{new}} \gets u + \alpha\,\nabla_u\mathcal{H}(t,z,p_x,u)$
        \If{$\|u_{\mathrm{new}}-u\| < \varepsilon_{\mathrm{tol}}$}
          \State $u \gets u_{\mathrm{new}}$
          \State \textbf{break}
        \EndIf
        \State $u \gets u_{\mathrm{new}}$
      \EndFor
      \For{$j = 1$ to $K'$} \Comment{$K'$ tracked steps; gradient flows through $\phi_\theta$}
        \State $u \gets u + \alpha\,\nabla_u\mathcal{H}(t,z,p_x,u)$
      \EndFor
      \State $u_k^\star \gets u$
      \State $\mathrm{running\_cost}_i \gets \mathrm{running\_cost}_i + h\,L(t,z,u_k^\star)$
      \State $z \gets z + h\,f(t,z,u_k^\star)$ \Comment{detach state transitions in code}
    \EndFor
    \State $\mathrm{terminal\_cost}_i \gets \alpha_G\,G(z)$
  \EndFor
  \State $\widetilde{\mathcal{L}}(\theta) \gets \frac{1}{B}\sum_{i=1}^{B}\bigl[\alpha_L\,\mathrm{running\_cost}_i + \mathrm{terminal\_cost}_i\bigr]$
  \State $\theta \gets \theta - \eta\,\nabla_\theta\widetilde{\mathcal{L}}(\theta)$ \Comment{JFB gradient (Term I)}
\Until{convergence}

\State \textbf{Deployment:} compute $p_x=\nabla_x\phi_{\theta^\star}(t,x)$ and iterate $u \gets u + \alpha\,\nabla_u\mathcal{H}(t,x,p_x,u)$ until convergence.
\end{algorithmic}
\end{algorithm}




\subsubsection{Computational Cost Per Iteration}

The dominant cost per iteration is $O(B \cdot M \cdot P)$ --- linear in batch
size $B$, time steps $M$, and parameters $P$. The fixed-point iterations add
$O(KBMm)$, which is negligible for $m \ll P$. This is the same cost as
training a standard neural network on $BM$ fixed data points.

Compare with exact backpropagation: the adjoint method adds
$O(B \cdot M \cdot m \cdot P)$ for the IFT term \eqref{eq:dustar_dtheta}.
For $n = 20$, $m = 2$, $P = 10^5$, JFB is approximately 20 times cheaper
per iteration.

\medskip\hrule\medskip

\subsubsection{Deployment: Using the Trained Network}

Once training is complete, the feedback law is deployed as follows. At any
state $x$ at time $t$:

\begin{enumerate}
    \item Compute $\nabla_x\phi_{\theta^\star}(t,x)$ by autodiff.
    \item Run fixed-point iterations $u \leftarrow u + \alpha\,
          \nabla_u\mathcal{H}(t,x,\nabla_x\phi_{\theta^\star},u)$ until convergence.
    \item Apply $u^\star$ to the system.
\end{enumerate}

The deployment cost per control evaluation is $O(P + Km)$ --- one network
forward-backward pass plus $K$ fixed-point steps. This is real-time capable for
moderate $P$ and small $m$, which covers the physical problems in the paper.
Crucially, deployment requires \textit{no} BVP solve, no trajectory
integration, and no knowledge of the initial condition. The entire optimal
policy for all initial conditions is encoded in $\theta^\star$.

\medskip

\hrule

\subsection{Synthesis: The Complete Contribution}
\label{sec:synthesis}

\medskip\hrule\medskip

\subsubsection{The Three Ingredients and Why Each Is Necessary}

The paper's method rests on three ingredients that must all be present
simultaneously. Removing any one of them breaks the method.

\textbf{Ingredient 1: Value function parameterization $\phi_\theta$.}
This is necessary rather than policy parameterization $u_\theta$ for three
reasons established in Section~\ref{sec:nn-param}: it injects physical structure via the HJB
constraint, it provides the costate $\nabla_x\phi_\theta$ needed for the
implicit solve, and it automatically satisfies the integrability condition
of the Lagrangian submanifold structure. Without $\phi_\theta$, there is no
costate, no Hamiltonian structure, and no implicit solve.

\textbf{Ingredient 2: The PMP--HJB bridge with implicit fixed-point solve.}
The identity $p_x = \nabla_x\phi_\theta$ converts the implicit optimality
condition $\nabla_u\mathcal{H}(t,z_x,p_x,u^\star) = 0$ into a fixed-point
equation solvable by gradient ascent on $\mathcal{H}$ w.r.t.\ $u$, without requiring a closed-form $g$.
This is what extends the class of problems from explicit to implicit
Hamiltonians. Without this ingredient, the method reduces to prior work that
requires $g$.

\textbf{Ingredient 3: Jacobian-Free Backpropagation.} The fixed-point solve makes
trajectory generation possible for implicit Hamiltonians, but it also adds
a layer of coupling between $\theta$ and the trajectories through
\eqref{eq:dustar_dtheta}. JFB removes this coupling by stopping gradients
through both the trajectory and the fixed-point iterations, reducing exact
$O(BmMP)$ backpropagation to $O(BMP)$ autodiff while preserving the descent
direction guarantee. Without JFB, the training loop is computationally
intractable for networks with more than a few thousand parameters.

\medskip

In one sentence: \textit{no prior method could learn a feedback controller
for high-dimensional optimal control problems with implicit Hamiltonians,
because all value-function methods required a closed-form
$u^\star = g(t,z_x,\nabla_x\phi)$ and all direct policy methods ignored
the Hamiltonian structure needed to handle the implicit case.}

\medskip\hrule\medskip

\subsubsection{Open Questions and Limitations}

For completeness, we note the assumptions on which the method rests and the
questions they leave open.

\textbf{Strict concavity of $\mathcal{H}$ in $u$.} This is required for
uniqueness of $u^\star$, invertibility of the Jacobian $\nabla^2_{uu}\mathcal{H}$, and the IFT
smoothness of $u^\star(\theta)$. For problems where $\mathcal{H}$ is not
concave in $u$ --- e.g., bang-bang problems where $\mathcal{H}$ is linear
in $u$ --- the fixed-point solve has no unique solution and the method as stated
does not apply. Extending to non-concave or discontinuous optimal controls
remains an open problem.

\textbf{Smoothness of $\phi$.} The bridge identity and the control
objective both require $\phi \in C^1$. When the value function develops kinks
(caustics), the network $\phi_\theta$ will attempt to approximate a
non-smooth function with a smooth one, introducing approximation error that
does not decrease with network size. Incorporating viscosity solution
structure into the architecture is an open research direction.

\textbf{Scalability of the fixed-point solve.} For control dimensions $m \gg 1$,
the cost of each gradient ascent step is $O(m)$, but more iterations may be
needed for convergence (larger contraction constant $\gamma$).
The paper demonstrates the method for $m = 1$ (space shuttle) and $m = 2$
(bicycle); scaling to $m \sim 10$ or higher remains to be demonstrated.

\end{document}